\documentclass[a4paper,reqno,11pt]{amsart}
\usepackage{a4wide}
%% \textheight 220mm
%% %\textwidth 160mm
%% %\textheight 230mm
%% \textwidth 150mm
%% \hoffset -16mm
%% %\voffset -16mm

%\usepackage{mathdots}
\usepackage{soul} %for strikethough
\usepackage[all]{xy}
%\usepackage{tixz-cd}
\usepackage{verbatim}
\usepackage{amstext}
\usepackage{amsmath}
\usepackage{amscd}
\usepackage{amsthm}
\usepackage{amsfonts}
%\usepackage{txfonts}
\usepackage{enumerate}
\usepackage{graphicx}
\usepackage{latexsym}
\usepackage{mathrsfs}
%\usepackage{showkeys}
%\usepackage[all,color,dvips]{xy}
\usepackage{mathtools}
%\input xy
\xyoption{all}

\usepackage{pstricks}
\usepackage{lscape}
%\usepackage[dvips]{color}
\usepackage{comment}

\newtheorem{theorem}{Theorem}[section]
\newtheorem*{theorem*}{Theorem}
\newtheorem{acknowledgement}[theorem]{Acknowledgement}
\newtheorem{corollary}[theorem]{Corollary}
\newtheorem{lemma}[theorem]{Lemma}
\newtheorem{proposition}[theorem]{Proposition}
\newtheorem{definition-proposition}[theorem]{Definition-Proposition}
\newtheorem{problem}[theorem]{Problem}
\newtheorem{question}[theorem]{Question}
\newtheorem{conjecture}[theorem]{Conjecture}

\theoremstyle{definition}
\newtheorem{definition}[theorem]{Definition}
\newtheorem{assumption}[theorem]{Assumption}
\newtheorem{remark}[theorem]{Remark}
\newtheorem{example}[theorem]{Example}
\newtheorem{observation}[theorem]{Observation}
\newtheorem{construction}[theorem]{Construction}

\newcommand{\mm}{{\mathfrak{m}}}
\newcommand{\nn}{{\mathfrak{n}}}
\newcommand{\pp}{{\mathfrak{p}}}
\newcommand{\qq}{{\mathfrak{q}}}
\renewcommand{\AA}{\mathscr{A}}
\newcommand{\CC}{\mathcal{C}}
\newcommand{\CCC}{\mathsf{C}}
\newcommand{\DDD}{\mathsf{D}}
\newcommand{\FF}{\mathscr{F}}
\newcommand{\KKK}{\mathsf{K}}
\newcommand{\PPP}{\mathcal{P}}
\newcommand{\PP}{\mathscr{P}}
\newcommand{\II}{\mathscr{I}}
\newcommand{\RR}{\mathscr{R}}
\renewcommand{\SS}{\mathscr{S}}
\newcommand{\TT}{\mathcal{T}}
\newcommand{\UU}{\mathcal{U}}
\newcommand{\MM}{\mathsf{M}}
\newcommand{\XX}{\mathscr{X}}
\newcommand{\YY}{\mathscr{Y}}
\newcommand{\w}{\vec{\omega}}
\renewcommand{\L}{\mathbb{L}}
\renewcommand{\P}{\mathbb{P}}
\newcommand{\Z}{\mathbb{Z}}
\newcommand{\R}{\mathbb{R}}
\newcommand{\C}{\mathbb{C}}
\newcommand{\X}{\mathbb{X}}
\renewcommand{\L}{\mathbb{L}}
\newcommand{\bo}{\operatorname{b}\nolimits}
\newcommand{\depth}{\operatorname{depth}\nolimits}
\newcommand{\projdim}{\operatorname{projdim}\nolimits}
\newcommand{\gldim}{\operatorname{gldim}\nolimits}
\newcommand{\soc}{\operatorname{soc}\nolimits}
\newcommand{\Ext}{\operatorname{Ext}\nolimits}
\newcommand{\Out}{\operatorname{Out}\nolimits}\newcommand{\Pic}{\operatorname{Pic}\nolimits}
\newcommand{\Tor}{\operatorname{Tor}\nolimits}
\newcommand{\Hom}{\operatorname{Hom}\nolimits}
\newcommand{\rad}{\operatorname{rad}\nolimits}
\newcommand{\End}{\operatorname{End}\nolimits}
\newcommand{\Aut}{\operatorname{Aut}\nolimits}
\newcommand{\gl}{\operatorname{gl.\!dim}\nolimits}
\newcommand{\injdim}{\operatorname{injdim}\nolimits}
\newcommand{\op}{\operatorname{op}\nolimits}
\newcommand{\sg}{\operatorname{sg}\nolimits}
\newcommand{\RHom}{\mathbf{R}\strut\kern-.2em\operatorname{Hom}\nolimits}
\newcommand{\Lotimes}{\mathop{\stackrel{\mathbf{L}}{\otimes}}\nolimits}
\newcommand{\Image}{\operatorname{Im}\nolimits}
\newcommand{\Kernel}{\operatorname{Ker}\nolimits}
\newcommand{\Cokernel}{\operatorname{Cok}\nolimits}
\newcommand{\Spec}{\operatorname{Spec}\nolimits}
\newcommand{\Supp}{\operatorname{Supp}\nolimits}
\newcommand{\MaxSpec}{\operatorname{MaxSpec}\nolimits}
\newcommand{\Ass}{\operatorname{Ass}\nolimits}
\newcommand{\ann}{\operatorname{ann}\nolimits}
\newcommand{\Proj}{\operatorname{Proj}\nolimits}
\newcommand{\MaxProj}{\operatorname{MaxProj}\nolimits}
\newcommand{\sign}{\operatorname{sgn}\nolimits}
\newcommand{\GL}{\operatorname{GL}\nolimits}
\newcommand{\SL}{\operatorname{SL}\nolimits}

\newcommand{\Aff}{\operatorname{Aff}\nolimits}

\newcommand{\Ho}{{\rm H}}
\newcommand{\LHo}{\underline{\rm H}}


\newcommand{\arr}[1]{\overset{#1}{\rightarrow}}
%\newcommand{\LM}[1]{\mathop{{#1}?}}
%\newcommand{\RM}[1]{\mathop{?{#1}}}
\newcommand{\LM}[1]{\mathop{L_{#1}}}
\newcommand{\RM}[1]{\mathop{R_{#1}}}
\newcommand{\gen}[1]{\langle #1 \rangle}
\renewcommand{\c}{\vec c}
\newcommand{\y}{\vec y}
\newcommand{\x}{\vec x}
\newcommand{\com}[1]{\operatorname{com}\nolimits_{#1}}

\DeclareMathOperator{\moduleCategory}{\mathsf{mod}} \renewcommand{\mod}{\moduleCategory}
\DeclareMathOperator{\Mod}{\mathsf{Mod}}
\DeclareMathOperator{\proj}{\mathsf{proj}}
\DeclareMathOperator{\inj}{\mathsf{inj}}
\DeclareMathOperator{\ind}{\mathsf{ind}}
\newcommand{\Db}{\mathrm{D}^{\mathrm{b}}}
\newcommand{\Kb}{\KKK^{\mathrm{b}}}
\DeclareMathOperator{\Sub}{\mathsf{Sub}}
\DeclareMathOperator{\Tors}{\mathsf{Tors}}
\DeclareMathOperator{\Torf}{\mathsf{Torf}}
\DeclareMathOperator{\thick}{\mathsf{thick}}
\DeclareMathOperator{\coh}{\mathsf{coh}}
\DeclareMathOperator{\per}{\mathsf{per}}
\DeclareMathOperator{\grproj}{\mathsf{gr-proj}}
\DeclareMathOperator{\Gr}{\mathsf{Gr}}
\DeclareMathOperator{\gr}{\mathsf{gr}}
\DeclareMathOperator{\qGr}{\mathsf{qGr}}
\DeclareMathOperator{\qgr}{\mathsf{qgr}}
\DeclareMathOperator{\CM}{\mathsf{CM}}
\DeclareMathOperator{\add}{\mathsf{add}}
\DeclareMathOperator{\fl}{\mathsf{f.\!l.}}
\DeclareMathOperator{\fd}{\mathsf{fd}}

\DeclareMathOperator{\repdim}{repdim}

%%% Added by Erik %%%
\DeclareMathOperator{\res}{res}
\newcommand{\stmod}{\mathop{\underline{\mod}}\nolimits}
\newcommand\kbpl{\Kb(\proj\Lambda)}


\newcommand{\new}[1]{{\blue #1}}
\newcommand{\old}[1]{{\red #1}}
\renewcommand{\comment}[1]{{\color{green!50!gray} #1}}

%\newcommand{\new}[1]{#1}
%\newcommand{\old}[1]{}
%\renewcommand{\comment}[1]{}

\newcommand{\cut}{\ar@{-}@[|(5)]}

\newcommand{\iso}{\cong}
\newcommand{\equi}{\approx}

\newcommand{\dimvec}{\underline{\operatorname{dim}}}
\newcommand{\sttilt}{\mathsf{s}\tau\mathsf{tilt}}
\newcommand{\tors}{\mathsf{tors}}
\DeclareMathOperator{\Fac}{\mathsf{Fac}}
\newcommand{\cyc}{\operatorname{cyc}\nolimits}
\newcommand{\roof}{\operatorname{roof}\nolimits}
\numberwithin{equation}{section}

\begin{document}
\title[Periodicity of trivial extension algebras]{Periodicity of trivial extension algebras}

\comment{Title is simplified. /OI}

\author[Chan]{Aaron Chan}

\author[Darp\"o]{Erik Darp\"o}

\author[Iyama]{Osamu Iyama}
\address{O. Iyama: Graduate School of Mathematics, Nagoya University, Furocho, Chikusaku, Nagoya 464-8602, Japan}
\email{iyama@math.nagoya-u.ac.jp}
\urladdr{http://www.math.nagoya-u.ac.jp/~iyama/}

\author[Marczinzik]{Rene Marczinzik}
\maketitle 

\begin{abstract}
\new{We study the periodicity (respectively, twisted periodicity) of the trivial extension algebra $T(A)$ of a finite dimensional algebra $A$. We prove that it is equivalent to the fractional Calabi-Yau property (respectively, twisted fractional Calabi-Yau property) of $A$. We also prove that twisted periodicity is equivalent to the $d$-representation-finiteness of the $r$-fold trivial extension algebra $T_r(A)$ for some $r,d\ge1$.}
\end{abstract}

\tableofcontents


\section*{Introduction}
The \emph{trivial extension algebra} of a finite dimensional algebra $A$ over a field $k$ is $T(A):=A\oplus DA$ with an obvious multiplication rule. This is one of the most basic classes of \emph{symmetric algebras}, that is, $B$ and its $k$-dual $DB$ are isomorphic as bimodules over itself.
The trivial extension algebra and its variation called the \emph{repetitive algebra} together with their dg analogues have played important roles in representation theory, for example, the classification of ($d$-)representation-finite or representation-tame self-injective algebras \cite{T,HW,R,Sk,SY,J}, the study of derived categories \cite{Hap} and the cluster categories \cite{K,Am}, the gentle algebras \cite{AS} and Brauer graph algebras \cite{Sc}, etc.

The first aim of this paper is to study the periodicity of trivial extension algebras and its application to cluster tilting theory. In the class of symmetric algebras (and self-injective algebras more generally), those satisfying the periodicity of the syzygy functor are considered to be the most basic ones (see Definition \ref{def:periodic} for the precise definitions and Propositions \ref{thm:hani}, \ref{thm:hani-gr} and \ref{gr=ungr} for known equivalent conditions).
For example, the trivial extension algebra of the path algebra $kQ$ of an acyclic quiver $Q$ is periodic if and only if $Q$ is Dynkin \cite{BBK}. In this case, $T(kQ)$ is $(2h-2)$-periodic for the Coxeter number $h$ of $Q$.
Our motivating questions are the following.

\begin{question}
%\begin{enumerate}[\rm(a)]
When is $T(A)$ periodic (or twisted periodic more generally)?
%\item When is $T_r(A)$ $d$-representation finite for some $d,r\ge1$?
%\end{enumerate}
\end{question}

We will give a complete answer in terms of the homological properties of $A$.  Moreover, as an application, we give a large number of new examples of periodic algebras \new{- including many representation-wild ones}.

For a finite dimensional algebra $A$ with finite global dimension, the bounded derived category $\DDD^{\bo}(\mod A)$ has a Serre functor $\nu=-\Lotimes_A(DA)$. We call $A$ \emph{fractionally Calabi-Yau} if there exist integers $\ell>0$, $m$ such that $\nu^\ell$ and $[m]$ are isomorphic as functors on $\DDD^{\bo}(\mod A)$. \marginpar{There is a stronger version: ``$\nu^\ell\simeq[m]$ as triangle functors'' [Keller, Calabi-Yau triangulated categories]} In this case, $A$ is called \old{(twisted)} $\frac{m}{\ell}$-Calabi-Yau.  For example, the path algebra of a Dynkin quiver is $\frac{h-2}{h}$-Calabi-Yau for the Coxeter number $h$.
There are many important examples of (twisted) fractionally Calabi-Yau algebras\new{; see, for example,} \cite{HerIya}.
There is also a weaker notion of \emph{twisted} fractionally Calabi-Yau algebras.  \comment{say something more?}

Our first main result characterizes the periodicity of $T(A)$ in terms of $A$.
%characterizes the periodicity of $A$ in terms of the fractionally Calabi-Yau property of $A$.

\begin{theorem}\label{main theorem}
Let $A$ be a finite dimensional $k$-algebra.
Then $T(A)$ is periodic if and only if $A$ has finite global dimension and is fractionally Calabi-Yau.
\end{theorem}

The trickiest part of the proof is the ``if'' part \old{for non-twisted case}, which will be shown in Section 3 by using the bar resolution of a certain differential graded algebra.

We also give several characterizations of the twisted periodicity of $T(A)$.
Recall that, for a positive integer $d$, a finite dimensional algebra $A$ is called \emph{$d$-representation finite} if there exists a $d$-cluster tilting $A$-module. This is closely related to the periodicity since $d$-representation finiteness implies a weaker version of the periodicity, that is, all modules have complexity at most one \cite{EH}. We study $d$-representation finiteness of $T_r(A)$ by applying the methods given in \cite[Section 2.2]{DI}.
Our second main result can be summarized as follows.

\begin{theorem}[=Theorem \ref{main theorem 2 again}]\label{main theorem 2}
Let $A$ be a finite dimensional algebra. Then the following conditions are equivalent.
\begin{enumerate}[\rm(i)] 
\item $T(A)$ is twisted periodic.
\item Each $T(A)$-module has complexity at most one.
\item There exist $d,r\ge1$ such that $T_r(A)$ is $d$-representation finite.
\item $A$ has finite global dimension and is twisted fractionally Calabi-Yau.
\end{enumerate}
\end{theorem}


%\begin{theorem}\label{main theorem}
%Let $A$ be a finite dimensional $k$-algebra.
%\begin{enumerate}[\rm(a)]
%\item The following conditions are equivalent.
%\begin{enumerate}[\rm(i)]
%\item $T(A)$ is periodic.
%\item $A$ has finite global dimension and is fractionally Calabi-Yau.
%\end{enumerate}
%\item The following conditions are equivalent.
%\begin{enumerate}[\rm(i)] 
%\item $T(A)$ is twisted periodic.
%\item There exists $d,r\ge1$ such that $T_r(A)$ is $d$-representation finite.
%\item $A$ has finite global dimension and is twisted fractionally Calabi-Yau.
%\end{enumerate}
%\end{enumerate}
%\end{theorem}


Theorem \ref{main theorem} gives the following implications.
\begin{align*}
\vcenter{
\xymatrix@C=60pt{
{\phantom{ed}A:\,\frac{m}{\ell}\text{-CY}\phantom{graddd}} \ar@{=>}[r]^{\text{trivial}\phantom{abcd}} \ar@{<=>}[d]^{\text{Thm }\ref{main theorem}} & {\quad\;\; A:\text{ twisted }\frac{m}{\ell}\text{-CY}\qquad} \ar@{<=>}[d]^{\text{Thm }\ref{main theorem 2}}\\
%B:\text{ graded periodic} \ar@{=>}[r]^{\text{trivial}\phantom{abcd}} \ar@{<=>}[d]_{\text{Prop }\ref{gr=ungr}}& B:\text{ graded twisted periodic}\ar@{<=>}[d]^{\text{Prop }\ref{gr=ungr}}\\
{\phantom{ed}T(A):\text{ periodic}\phantom{grad}}  \ar@{=>}[r]^{\text{trivial}\phantom{abcd}} & {\phantom{ged}T(A):\text{ twisted periodic}\phantom{rad}}
}}
\end{align*}
There are also other natural variations of the periodicty, see Section 2. We will make clear the connection between those conditions, see \eqref{dgm:implications}. On the other hand, the following open question asks if all the conditions in the diagram above are equivalent.

\begin{question}\label{conj:periodicity}
\cite{ES} A finite dimensional algebra is periodic if and only if it is twisted periodic.
\end{question}

One of the consequences of Theorem \ref{main theorem} is the following.

\begin{corollary}
Let $A$ be a finite dimensional algebra with finite global dimension. If the outer automorphism group of $A$ is finite, then Question \ref{conj:periodicity} is true for $T(A)$.
\end{corollary}

For example, Question \ref{conj:periodicity} is true for the trivial extension algebra of the incidence algebra $A$ of a lattice $P$ (or a bounded poset more generally) since the outer automorphism group of $A$ is canonically isomorphic to the automorphism group of $P$ which is clearly finite.

There is a similar type of question for fractionally Calabi-Yau algebras.

\begin{question}\cite{HerIya}\label{remove twist 2}
A finite dimensional algebra of finite global dimension is fractionally Calabi-Yau if and only if it is twisted fractionally Calabi-Yau.
\end{question}

Theorem \ref{main theorem} shows that these two questions are equivalent in the following sense.

\begin{corollary}
Let $A$ be a a finite dimensional algebra of finite global dimension. Then Question \ref{conj:periodicity} is true for $T(A)$ if and only if Question \ref{remove twist 2} is true for $T(A)$.\end{corollary}






Notice that the $d$-representation finiteness of $T_r(A)$ implies the $d$-representation finiteness of $T_{r\ell}(A)$ for all $\ell\ge1$. In Section 8, we give a family of $A$ such that $T(A)$ is $d$-representation finite, which is the strongest condition in the sense above. \comment{State a main result in Section 4 here.}

\comment{Also explain results and examples from the incidence algebras:}

\begin{example}
Let $A=(k\mathbb{A}_2)^{\otimes n}$. Then $T(A)$ is ...
\end{example}


\comment{The following should be more seriously considered and written somewhere else.}
It is natural to pose the following question.

\begin{question}
Is the order of the outer automorphism group of each (twisted)  fractionally Calabi-Yau algebra finite?
\end{question}

Notice that the notion of (twisted) fractionally Calabi-Yau algebras can be defined for Gorenstein algebras more generally. Any self-injective algebra is twisted $\frac{0}{1}$-Calabi-Yau whose twist is given by the Nakayama automorphism. But they are usually not fractionally Calabi-Yau since the Nakayama automorphism usually have infinite order.


\section{Preliminaries}

Unless otherwise specified, by $\Lambda$-module we mean finitely generated right $\Lambda$-module.
The category of $\Lambda$-modules are denoted by $\mod \Lambda$.
Denote by $D:=\Hom_k(-,k)$ the $k$-linear duality.
\marginpar{We should fix our convention for $\phi^*$ and $\phi_*$. /OI}
\new{For an algebra automorphism $\phi$ of $A$, we denote by $\phi_*$ the induced automorphism of $\mod A$.

Let $\Pic_k(A)$ be the Picard group of $A$. Thus the element of $\Pic_k(A)$ is an isomorphism class of an $A^e$-module $X$ such that there exists an $A^e$-module $Y$ such that $X\otimes_AY\simeq A\simeq Y\otimes_X$ as $A^e$-modules, and the multiplication is given by tensor product.
If $A$ is basic, then there is a canonical isomorphism
\[\Out_k(A)\simeq\Pic_k(A)\ \mbox{ given by }\ \alpha\mapsto{}_1A_\alpha.\]
}\marginpar{This is frequently used. /OI}


\subsection{Trivial extension algebras}
Recall that the \emph{trivial extension} $T(A)$ of a finite dimensional algebra $A$ is by definition $T(A)=A \oplus D(A)$ as vector spaces with multiplication $(a,f)(b,g)=(ab,ag+fb)$ for $a,b \in A$ and $f,g \in D(A)$.

The trivial extension algebra is an important tool in ring theory and homological algebra and we refer for example to the textbook \cite{FGR} for more on trivial extensions in a more general context.
We collect some elementary properties and refer for example to \cite{SY} for proofs.
\begin{proposition} \label{propotrivext}
Let $A$ be a finite dimensional algebra.
\begin{enumerate}
\item $T(A)$ is a symmetric Frobenius algebra.
\item $\dim(T(A))=2 \dim(A)$.
\item $\rad(T(A))=(\rad(A),D(A))$.
\end{enumerate}

\end{proposition}

\subsection{Serre functor and \new{fractionally Calabi--Yau algebras}}\label{subsec:CYdim}

Suppose $A$ is a finite-dimensional algebra of finite global dimension.
Then we have an auto-equivalence $\nu$ on the bounded derived category $\Db(\mod A)$ of $A$-modules given by
\[
\nu:= \Lotimes_ADA = D\circ \RHom_A(-,A).
\]
This is the \emph{Serre functor} on $\Db(\mod A)$, meaning that we have a bifunctorial isomorphism $\Hom_{\Db(\mod A)}(X,Y)\cong D\Hom_{\Db(\mod A)}(Y,\nu(X))$.


For a finite-dimensional algebra $A$ of finite global dimension, and the trivial extension $B:=T(A)$ of $A$, we have Happel's triangle equivalence \cite{Hap}:\marginpar{These diagrams should come first. /OI}
\begin{equation}\label{happel}
\DDD^{\bo}(\mod A)\simeq\stmod^{\Z} B.
\end{equation}
By uniqueness of the Serre functor, we have a commutative diagram
\begin{equation}\label{serre}
\xymatrix{
\DDD^{\bo}(\mod A)\ar[r]^{\sim}\ar[d]^\nu&\stmod^{\Z}B\ar[d]^{\Omega(1)}\\
\DDD^{\bo}(\mod A)\ar[r]^{\sim}&\stmod^{\Z}B.}
\end{equation}



\begin{definition}
Suppose $A$ is a finite-dimensional algebra of finite global dimension.
For non-zero integers $\ell$ and $m$, $A$ is said to be \emph{twisted $\frac{m}{\ell}$-Calabi--Yau} if there is an automorphism $\phi$ of $A$ so 
we have a natural isomorphism of auto-equivalences
\begin{equation}\label{eq:tfCY}
\nu^\ell \simeq [m]\circ \phi^*
\end{equation}
on $\Db(\mod A)$.
Note that this is equivalent to $\nu^\ell(A)\cong A[m]$ in $\Db(\mod A)$; see \cite[Prop 4.3]{HerIya}.
If, moreover, $\phi=\mathrm{id}$, we say that $A$ is \emph{$\frac{m}{\ell}$-Calabi--Yau}.
We say that $\frac{m}{\ell}$ is the \emph{minimal Calabi--Yau dimension} of $A$ if $\ell$ is the smallest positive integer such that \eqref{eq:tfCY} holds.
\end{definition}

\marginpar{TODO: comment about field extension\\
About what? /OI}

\begin{example}\label{eg:fCYdim}
(1) The path algebra $kQ$ associated to a quiver $Q$ of Dynkin type $\Delta$ is $\frac{h-2}{h}$-Calabi--Yau, where $h$ is the Coxeter number of $\Delta$.  %\new{The minimal Calabi--Yau dimension is $\frac{h-2}{h}$ if the orientation of $Q$ is not alternating; otherwise, it is $\frac{h/2-1}{h/2}$.}

(2) \cite[Proposition 1.4]{HerIya}
%\marginpar{needs $A_i/\rad(A_i)$ separable. /ED}  
Suppose $A_i$ is \new{finite dimensional $k$-algebra with finite global dimension which is $\frac{m_i}{l_i}$-Calabi--Yau for $i\in\{1,\ldots, t\}$. If $A_1\otimes_kA_2$ has finite global dimension, then the algebra $A= \bigotimes\limits_{i=1}^{t}{\Lambda_i}$ is $\frac{m}{l}$-Calabi--Yau} with 
\[
l=\mathrm{lcm}(l_1,...,l_t) \text{ and } m=l\left(\frac{m_1}{l_1}+\cdots+ \frac{m_t}{l_t}\right).
\]

\new{(3) Let $kQ$ be the path algebra of Dynkin quiver and $\Lambda$ the stable Auslander algebra of $kQ$. Then it is fractionally Calabi-Yau.\marginpar{I will write more details later. /OI}

(4) Let $A$ be a canonical algebra of tubular type, that is, $(2,2,2,2)$, $(3,3,3)$, $(2,4,4)$ or $(2,3,6)$. Then $A$ is $\frac{p}{p}$-Calabi-Yau for $p=2,3,4$ or $6$.

More generally, let $A$ be a $d$-canonical algebra of type $(p_1,\ldots,p_n)$ such that $n-d-1=\sum_{i=1}^n\frac{1}{p_i}$. Then $A$ is $\frac{p}{p}$-Calabi-Yau for $p={\rm l.c.m}(p_1,\ldots,p_n)$ \cite{HIMO}.
}

\end{example}


\subsection{Cluster tilting}

\begin{definition}
Let $\mathcal{E}$ be an exact category and $\mathcal{T} \subseteq \mathcal{E}$ a full and functorially finite subcategory closed under direct sums and direct summands. Then $\mathcal{T}$ is called an \emph{$m$-cluster-tilting subcategory} if 
$$\{ X \in \mathcal{E} \mid \Ext_{\mathcal{E}}^i(X,\mathcal{T})=0, \ 0 < i < m \}=\mathcal{T}=\{ X \in \mathcal{E} \mid \Ext_{\mathcal{E}}^i(\mathcal{T},X)=0, \ 0 < i < m \}.$$
\end{definition}

\old{For the next definition we refer for example to section 3 in \cite{P}.
\begin{definition}
A Frobenius category $\mathcal{E}$ is called \emph{stably $m$-Calabi--Yau} for some $m \geq 1$ if its stable category $\underline{\mathcal{E}}$ is $m$-Calabi--Yau, which means that $\underline{\mathcal{E}}$ is Hom-finite and there is a functorial duality
$$D \underline{\Hom}_{\mathcal{E}}(X,Y) = \underline{\Hom}_{\mathcal{E}}(Y,\Omega^{-m}(X)).$$
$\mathcal{E}$ is called a \emph{Frobenius $m$-cluster category} if it is idempotent complete, stably $m$-Calabi--Yau and $\gldim \End_{\mathcal{E}}(T) \leq m+1$ for any $m$-cluster tilting object $T$, of which there is at least one. 
\end{definition}

We say for short that a finite dimensional algebra $A$ is \emph{Frobenius $m$-cluster} in case $\mod A$ is a Frobenius $m$-cluster category.}
\marginpar{The conditions are fine, but the name is strange. /OI}

One of the most famous Frobenius $m$-cluster algebras are the preprojective algebras of Dynkin type, where we have $m=2$.
Since we have $D \underline{\Hom}_A(M,N) \cong \underline{\Hom}_A(N,\Omega^1(M))$ for all $M,N \in \mod A$ for a symmetric algebra $A$, the module category of $A$ is stably $m$-Calabi--Yau when $A$ is $m+1$-periodic, that is $\Omega_{A^e}^{m+1}(A) \cong A$ for the enveloping algebra $A^e$ of $A$.
One of the main applications of our results will be to show that trivial extension of the indicence algebra of the Boolean lattice on an $n$-set is a Frobenius $m$-cluster algebra for a suitable $m$ when $n$ is odd or the characteristic is two.

\new{EDIT Note (AC): I put the material that are about Boolean algebra as comment using the iffalse-fi command; they should be rearranged to, perhaps(?), a separation section in a later edition.}

\iffalse
\section{Problems and observations}
Let $A=\bigotimes_{i=1}^nk\mathbb{A}_2$ and $B=T(A)$.

Using QPA, Rene found MANY $(n+2)$-cluster tilting $A$-modules.  We can prove one (up to $\Omega$-shifts) of them is $(n+2)$-precluster tilting.  Notation for indices of simple $A$-modules are of the form $\pmb{x}=(x_1, \ldots, x_n)\in (\Z/2\Z)^n$, we will use the notation $\pmb{1}:=(1,1,\ldots, 1)$ and $\pmb{0}:=(0,0,\ldots, 0)$, and $|\pmb{x}|$ to denote the number of $1$'s in $\pmb{x}$ (i.e. the size of the subset of $\{1,2,\ldots, n\}$ corresponding to $\pmb{x}$).  Then our precluster tilting module is
\[
 B\oplus \bigoplus_{\pmb{x}\,:\, x_1=0} \Omega_B^{|\pmb{x}|}(S_{\pmb{x}})
\]
For an algebra $\Lambda$, if $X\in\mod\Lambda$ is \emph{weakly $\nu$-formal} in the sense that $\nu^i(X)$ is a stalk complex for each $i\ge0$, then one can describe the projective resolution of $X$ as a $T(\Lambda)$-module explicitly using the method of \cite{CIM}.
%Strictly speaking Serre-formal module should be $\nu^i(X)$ being stalk for all $i\neq 0$.  But for this implication about projective resolution, $i\geq 0$ is enough.  Probably $i\geq 0$ case implies $i\neq 0$ case, but we never figure out how.

This can be used to write down explicitly the projective resolution (over both $A$ and $B$) of $S_{\pmb{x}}$.  In particular, we have $\Omega_B^{-(|\pmb{x}|+1)}(S_{\pmb{x}}) = D(Ae_{\pmb{1}-\pmb{x}})$ and $\Omega_B^{n-|\pmb{x}|+1}(S_{\pmb{x}}) = e_{\pmb{1}-\pmb{x}}A$.  Therefore, the above module can be $\Omega_B$-shifted to
\[
 B\oplus \bigoplus_{\pmb{x}\,:\,x_1=1} D(Ae_{\pmb{x}}) \quad\text{ and also } \quad  B\oplus \bigoplus_{\pmb{x}\,:\,x_1=1} e_{\pmb{x}}A.
\]

\begin{problem}\label{main problem}
\begin{enumerate}[\rm(a)]
\item Prove that $B$ has an $(n+2)$-cluster tilting module.
\item Prove that $\Omega^{n+3}\simeq{\rm id}$ as functors on $\stmod B$.
\item Generalize (a) and (b) to a more general class of algebras.
\item Let $\Lambda$ be a finite-dimensional algebra of finite global dimension. Are weakly $\nu$-formal $\Lambda$-modules are $\nu$-formal (i.e.\ $\nu^i(X)$ is a stalk complex for all $i\in\Z$)? This is clear if $\Lambda$ is twisted fractionally Calabi--Yau.
\end{enumerate}
\end{problem}

Regard $B$ as a $\Z$-graded algebra with $B_0=A$ and $B_1=DA$. We denote by $(1)$ the degree shift functor in $\mod^{\Z}B$.
Thanks to \cite[2.14(a)]{DI}, the above (a) is equivalent to the following statement.
\begin{enumerate}[\rm(a)]
\item[\rm(a')] Prove that $\mod^{\Z}B$ has a $(1)$-stable $(n+2)$-cluster subcategory $\CC$  such that $(\ind\CC)/(1)$ is a finite set.
\end{enumerate}
\fi


%%%%%%%%%%%%%%%%%%%%%%%%%%%%%%%%%%%%%%%%%%%%%%%%%%%%%%%%%%%%%%%%%%%
%%%%%%%%%%%%%%%%%%%%%%%%%%%%%%%%%%%%%%%%%%%%%%%%%%%%%%%%%%%%%%%%%%%
%%%%%%%%%%%%%%%%%%%%%%%%%%%%%%%%%%%%%%%%%%%%%%%%%%%%%%%%%%%%%%%%%%%
%%%%%%%%%%%%%%%%%%%%%%%%%%%%%%%%%%%%%%%%%%%%%%%%%%%%%%%%%%%%%%%%%%%


\section{\new{Preliminaries on periodicity and twisted periodicity}}%{Interaction between periodicity and Calabi--Yau dimension}

Let us review the various notions of periodicity of algebras and modules.

\begin{definition} \label{def:periodic}
Suppose $\Lambda$ is a finite-dimensional algebra.
\begin{itemize}
\item $M\in \mod\Lambda$ is \emph{$\Omega$-periodic} if there is some integer $n\geq 0$ such that $\Omega^n(M)\cong M$ in $\stmod \Lambda$.  The number $n$ is the \emph{$\Omega$-period} of $M$ and is said to be \emph{minimal} if there is no $0\leq n'<n$ such that $\Omega^{n'}(M)\cong M$ in $\stmod\Lambda$.

%\item $\Lambda$ is \emph{functorially periodic with twist $\phi\in\Aut(\Lambda)$} if there is a natural isomorphism $\Omega^n\simeq (-)_\phi$ of autoequivalences on $\stmod \Lambda$ for some integer $n\geq 0$.  The number $n$ here is called the \emph{twisted functorial period} of $\Lambda$, and it is said to be \emph{minimal} if there is no $n'<n$ such that $\Lambda$ is functorially periodic of period $n'$.

\item $\Lambda$ is \emph{twisted (bimodule) periodic} if $\Omega_{\Lambda^e}^n(\Lambda)\cong {}_1\Lambda_\phi$ in $\stmod \Lambda^e$ for some integer $n\geq 0$ and $\phi\in \Aut(A)$.  We call the integer $n$ the \emph{twisted period} and $\phi$ the \emph{associated twist}.

\item $\Lambda$ is \emph{(bimodule) periodic} if it is twisted periodic with trivial associated twist $\phi=1$, i.e. $\Lambda$ is $\Omega$-periodic as an $\Lambda$-bimodule.
\end{itemize}
%\marginpar{Add $\Omega$-periodicity of $\stmod \Lambda$?}
\end{definition}

Clearly, periodic algebras are twisted periodic, but it is still open whether the converse holds, see Question \ref{conj:periodicity}. \marginpar{Question was moved to introduction. /OI}

%\begin{conjecture}[Periodicity conjecture]\label{conj:periodicity}
%If $\Lambda$ is a twisted periodic algebra, then it is periodic.
%\end{conjecture}


For graded algebras, all notions of periodicity admit a graded analogue, where any isomorphism appearing in the definition is taken up to a graded shift.\marginpar{\rule{0pt}{2.8ex}Write out the graded definitions}

Among self-injective algebras, trivial extension algebras form a very well-behaved subclass.  Note that trivial extension algebras have a canonical graded structure, given by $\deg A=0$ and $\deg DA=1$.


Let us start with reminding some equivalence conditions of twisted periodicity.

\begin{proposition}{\rm \cite{GSS, H}}\label{thm:hani}
For a ring-indecomposable finite-dimensional algebra $\Lambda$ with $\Lambda/\rad\Lambda$ separable over $k$, the following statements are equivalent.\marginpar{Written as Proposition. /OI}
\begin{enumerate}
\item All simple $\Lambda$-modules are $\Omega$-periodic.
\item $\Lambda$ is self-injective and for some $n\geq 0$ and some $\phi\in \Aut(\Lambda)$, there is a natural isomorphism $\Omega^n\simeq \phi_*$ of autoequivalences of $\stmod \Lambda$.\old{, where $\phi_*$ denote the natural functor induced by $\phi$.} \marginpar{$\phi_*$ is written above. /OI}
\item $\Lambda$ is twisted periodic.
\end{enumerate}
\end{proposition}

\new{We also have the following equivalent conditions of graded twisted periodicity.}
\old{A similar version of Proposition \ref{thm:hani} holds in the graded setting.}

\begin{proposition}{\rm \cite[2.4]{H}}\label{thm:hani-gr}
Let $\Lambda$ be a ring-indecomposable finite-dimensional $\Z$-graded algebra with $\Lambda/\rad\Lambda$ separable over $k$. Then the following statements are equivalent.\marginpar{Written as Proposition. /OI}
\begin{enumerate}
\item All simple $\Lambda$-modules are graded $\Omega$-periodic, i.e. for any simple $S$, there exist $n\in\Z_{>0}$ and $a\in\Z$ such that $\Omega^n(S)\cong S(a)$.

\item There exist integers $a,n$, and an idempotent-fixing graded automorphism $\phi$ of $\Lambda$ such that
  $\Omega^n\simeq \phi_*\circ(a)$ as automorphisms of $\stmod \Lambda$.
  %% For some idempotent-fixing graded automorphism $\phi$ of $\Lambda$, there is a natural isomorphism $\Omega^n\simeq \phi_*\circ(a)$ of autoequivalences on $\stmod \Lambda$ for some integers $n\geq 0$ and $a$.

\item $\Lambda$ is graded twisted periodic with an idempotent-fixing graded twist $\phi$, i.e. $\Omega_{\Lambda^e}^N(\Lambda)\cong {}_1\Lambda_\phi(a)$ in $\stmod^\Z \Lambda$ for some integers $a\in\Z$, $N\geq 0$.
\end{enumerate}
\end{proposition}

\old{\begin{remark}
Even if (2) holds with $\phi=1$, $\Lambda$ is not necessarily bimodule periodic. \marginpar{We do not need this remark. /OI}
\end{remark}}

Thankfully, the graded analogues are equivalent to the ungraded notion.  

\begin{proposition}\label{gr=ungr}
Let $\Lambda$ be a ring-indecomposable finite-dimensional $\Z$-graded algebra with $\Lambda/\rad\Lambda$ separable over $k$. The following statements are equivalent.
\begin{enumerate}[(1)]
\item $\Lambda$ is twisted periodic if, and only if, $\Lambda$ is graded twisted periodic with an idempotent-fixing graded twist.

\item $\Lambda$ is periodic if, and only if, $\Lambda$ is graded periodic, i.e. $\Omega_{\Lambda^e}(\Lambda)\cong \Lambda(a)$ in $\stmod^\Z \Lambda^e$ for some $a\in \Z$ and some $n>0$, with respect to the canonical graded structure on $\Lambda^e$ induced by that of $\Lambda$.
\end{enumerate}
\end{proposition}
\begin{proof}
For gradable $\Lambda$-modules $P,M$ with graded lifts $\widetilde{P}, \widetilde{M}$ respectively, it follows from the definition of $\mod^\Z\Lambda$ that we have
\[\Hom_{\Lambda}(P,M) = \Hom_{\Lambda}^\Z( \bigoplus_{a\in \Z}\widetilde{P}(a), \widetilde{M}). \] 
Since any projective module $P$ is gradable, if $P$ is the projective cover of $M$ in $\mod\Lambda$ with indecomposable decomposition $P=\bigoplus_{i\in I} P_i$, then the projective cover of $\widetilde{M}$ is $\bigoplus_{i\in I} P_i(a_i)$ for some integers $a_i\in \Z$.  Consequently, if the syzygy $\Omega(M)$ of $M$ in $\mod \Lambda$ has indecomposable decomposition $\Omega(M)=\bigoplus_{i\in I} N_i$, then the graded module $\Omega(\widetilde{M})$ is given by $\bigoplus_{i\in I}N_i(a_i)$ for some integers $a_i\in \Z$.

(1)  Note that every simple $\Lambda$-module can be graded as one that is concentrated in degree 0.  Hence, by Propositons \ref{thm:hani} and \ref{thm:hani-gr}, it suffices to show that a simple $\Lambda$-module $S$ is $\Omega$-periodic if and only if it is graded $\Omega$-periodic.  But this follows from the fact that syzygy of a gradable module coincides in the graded and ungraded setting up to shifts.

(2)  Follows immediate from the fact that syzygy of a gradable module coincides with the syzygy in the ungraded category up to shifts.
\end{proof}

\section{\new{Our results on periodicity and twisted periodicity}}
In this section, we shall give some relations between the different notions of periodicity. Our results are summarised in the diagram below. The implication indicated by a dotted arrow will be proved in Section~\ref{sec:fCY to per}.\marginpar{Extend \eqref{dgm:implications} by adding conditions in \ref{main theorem 2 again}? /OI}
\[\text{Assumptions: }\begin{cases}\text{$A$ is finite-dimensional with $\gl A<\infty$,}\\\text{$B:=T(A)$ is the trivial extension of $A$.}\end{cases}\]
\begin{align}\label{dgm:implications}
\vcenter{
\xymatrix@C=60pt{
{\phantom{ed}A:\,\frac{m}{\ell}\text{-CY}\phantom{graddd}} \ar@{=>}[r]^{\text{trivial}\phantom{abcd}} \ar@{:>}[d] & {\quad\;\; A:\text{ twisted }\frac{m}{\ell}\text{-CY}\qquad} \ar@{=>}[d]^{\text{Prop }\ref{tw-fCY to tw-pe}}\\
B:\text{ graded periodic} \ar@{=>}[r]^{\text{trivial}\phantom{abcd}} \ar@{<=>}[d]_{\text{Prop }\ref{gr=ungr} (2)} \ar@<3ex>@{=>}[u]^{\text{Prop }\ref{periodic to fCY}}& B:\text{ graded twisted periodic}\ar@{<=>}[d]^{\text{Prop }\ref{gr=ungr} (1)}\\
{\phantom{ed}B:\text{ periodic}\phantom{grad}}  \ar@{=>}[r]^{\text{trivial}\phantom{abcd}} & {\phantom{ged}B:\text{ twisted periodic}\phantom{rad}}
}}
\end{align}

\new{
In this Section, we prove the following result, which implies the right part of \eqref{dgm:implications}.

\begin{theorem}\label{main theorem 2 again}
Let $A$ be a finite dimensional algebra and $B=T(A)$. Then the following conditions are equivalent.
\begin{enumerate}[\rm(i)]
\item $B$ is twisted periodic.
\item $B$ is graded twisted periodic.
\item Each $B$-module has complexity at most one.
\item There exist $d,r\ge1$ such that $T_r(A)$ is $d$-representation finite.
\item $A$ has finite global dimension and is twisted fractionally Calabi-Yau.
\end{enumerate}
\end{theorem}
}

In the Section \ref{sec:fCY to per}, we will show that the dotted implication in the top left of the diagram holds.
If the Periodicity Conjecture \ref{conj:periodicity} holds (i.e. the implication arrow $\Rightarrow$ in the bottom row of \eqref{dgm:implications} can be upgraded to an equivalence $\Leftrightarrow$),
then all the conditions in the diagram are equivalent. In particular, every twisted fractionally Calabi--Yau algebra is also \emph{untwisted} fractionally Calabi--Yau (albeit possibly with different Calabi--Yau dimension)
  an algebra with twisted fractionally Calabi--Yau dimension necessarily has \emph{untwisted} fractionally Calabi--Yau dimension
  -- resolving a question in \cite[Remark 1.6(b)]{HerIya}.

\new{
\subsection{Application of Happel's equivalence}

Our main tool of the proof of Theorem \ref{main theorem 2 again} is Happel's equivalence \eqref{happel} and the commutative diagram \eqref{serre}.}
%Let us now come back to the setting of \eqref{dgm:implications}.  

%Since $k\mathbb{A}_2$ is $\frac{1}{3}$-Calabi--Yau, our $A$ is $\frac{n}{3}$-Calabi--Yau.
%\begin{equation}\label{n/3}
%\nu^3\simeq[n]
%\end{equation}
%as functors on $\DDD^{\bo}(\mod A)$.

%\new{\texttt{Insert here: Happel's equiv gives relation between periodicity of $B$ and CY properties of $A$.}}
\begin{proposition}\label{tw-fCY to tw-pe}
Let $A$ be a finite-dimensional algebra of finite global dimension, and $B:=T(A)$ be its trivial extension.
If $A$ is twisted $\frac{m}{\ell}$-Calabi--Yau, then \new{we have an isomorphism of functors
\[\Omega^{m+\ell}\simeq\phi_*\circ(-\ell)\ \mbox{ on }\ \stmod^\Z B.\]}
\old{every simple $B$-module $S$ satisfies $\Omega^{m+\ell}(S)\simeq S(-\ell)$ in $\stmod^\Z B$.}
In particular, such a $B$ is (graded) twisted periodic.
\end{proposition}

\begin{proof}
By \eqref{serre}, we have $\phi_*\simeq\nu^\ell[-m]\simeq \Omega^{\ell+m}(\ell)$ for some automorphism $\phi\in \Aut_k(A)$ of $A$.
\old{Since $\phi(1)=1$, we have $\Omega^{\ell+m}(A/\rad A)(\ell)=\phi_*(A/\rad A)=A/\rad A$ as required.}
It follows from Proposition \ref{thm:hani-gr} that in such a case, $B$ is graded twisted periodic, but this means that $B$ is twisted periodic, by Proposition \ref{gr=ungr}.
\end{proof}

If we consider the untwisted version, then we can obtain a if-and-only-if statement instead.
\begin{proposition}\label{periodic to fCY}
Let $A$ be a finite-dimensional algebra of finite global dimension, and $B:=T(A)$ its trivial extension. Given $m,\ell\in\Z$ with $\ell>0$, the following are equivalent.
\begin{enumerate}[(1)]
\item
  The algebra $A$ is $\frac{m}{\ell}$-Calabi--Yau.
\item
There is an isomorphism of functors $\Omega^{\ell+m}\simeq  (-\ell)$ on $\stmod^{\Z}B$.
\end{enumerate}
In particular, if $B$ is (graded) periodic, then $A$ is necessarily fractionally Calabi--Yau.
\end{proposition}

\begin{proof}
In verbatim to the proof of Proposition \ref{tw-fCY to tw-pe} by taking the identity automorphism $\phi=\mathrm{id}$ - that is, we have natural isomorphisms
$\phi_*=\mathrm{Id}_{\stmod^\Z B}=\nu^\ell[-m]= \Omega^{\ell+m}(\ell)$.  The claimed equivalence of conditions immediately follows.

For the final part, first note that, by Proposition \ref{gr=ungr}, $B$ being periodic implies that we have $\Omega_{B^e}^{\ell+m}(B)\cong B(-\ell)$ for some integers $\ell,m$.  Therefore, we have isomorphism of autoequvialences \[
\Omega_B^{\ell+m} = -\otimes_B\Omega_{B^e}^{\ell+m}(B) = -\otimes_BB(-\ell) = (-\ell)
\]
on $\stmod^\Z B$, i.e. the condition (2) holds.  Thus, it follows that $A$ is necessarily $\frac{m}{\ell}$-Calabi--Yau.
\end{proof}

\new{Next we show that the periodicity of $T(A)$ implies the finiteness of global dimension of $A$.} 
%We start by explaining why the assumption of $A$ having finite global dimension is natural when studying the periodicity of its trivial extension algebra $T(A)$. In particular, along with Proposition \ref{gr=ungr} and Happel's equivalence \eqref{happel}, this assumption will open up more tools for our disposal.

\begin{proposition}\label{fin gl}
Suppose $A$ is a finite-dimensional algebra and $B:=T(A)$ is its trivial extension.
If $B$ is (graded) twisted periodic, then $\gl A$ is finite.
\end{proposition}

\begin{proof}
A simple $A$-module $S$ is naturally a simple $B$-module concentrated in degree $0$.  By the periodicity assumption and Proposition~\ref{thm:hani}, we have an exact sequence
\[
0 \to S(i) \to P_n \to \cdots \to P_0 \rightarrow S\to 0
\]
in $\mod^{\Z}B$,
for some $i\in \Z_{\leq 0}$ with $P_j\in \proj^\Z B$ for all $j$.
\new{Then $i\le -1$ holds since each $P_i$ is generated in non-negative degrees and $S(i)\subset\soc P_n$.}
%First note that the grading shift $i$ cannot be zero because $S$ is a direct summand of $\soc(P_n)$ and, by the definition of grading on $B=T(A)$, $\soc(P) \simeq (P/\rad P)(-1)$ holds for every $P\in\proj^\Z B$.
%\old{the projective $P_n$ is supported in at least two (non-negative) degrees.}
Taking the degree $0$ part of the above exact sequence yields an exact sequence of $A$-modules
\[
0 \to (P_n)_0 \to \cdots \to (P_0)_0 \rightarrow S\to 0
\]
with $(P_j)_0$ projective $A$-module for all $j\geq 0$. Hence, the projective dimension of $S$ as an $A$-module is finite, and so is the global dimension of $A$.
\end{proof}

\new{
More technical argument gives the following stronger assertion.

\begin{proposition}\label{complexity one}
Let $A$ be a ring-indecomposable finite dimensional algebra, and $B:=T(A)$ its trivial extension. If each $B$-module has complexity at most one, then $A$ has finite global dimension and is twisted fractionally Calabi-Yau.
\end{proposition}

We need the following well-known result.

\begin{proposition}[e.g.\ \cite{HS}]\label{voigt}
Let $\Lambda$ be a finite dimensional algebra over a field $k$. For each $n\ge1$, there are only finitely many isomorphism classes of $\Lambda$-modules $X$ satisfying $\dim_kX=n$ and $\Ext^1_\Lambda(X,X)=0$.
\end{proposition}

\begin{proof}
When $k$ is algebraically closed, the assertion follows from Voigt's Lemma \cite{V}. For general case, let $\overline{k}$ be the algebraic closure of $k$ and $\overline{\Lambda}:=\overline{k}\otimes_k\Lambda$. The functor $\overline{(-)}:\mod\Lambda\to\mod\overline{\Lambda}$ satisfies $\dim_{\overline{k}}\overline{X}=\dim_kX$ and $\Ext^1_{\overline{\Lambda}}(\overline{X},\overline{X})=\overline{k}\otimes_k\Ext^1_\Lambda(X,X)$.
Moreover, $X\simeq Y$ holds as $\Lambda$-modules if and only if $\overline{X}\simeq\overline{Y}$ as $\overline{\Lambda}$-modules. Thus the assertion follows.
\end{proof}

We also need the following partial answer to the finitistic dimension conjecture.

\begin{proposition}\cite{JL,HS}\label{FDC}
For each $n,n'\ge1$, there exists an integer $m\ge1$ satisfying the following:
For each field $k$, a finite dimensional $k$-algebra $\Lambda$ with $\dim_kA\le n$ and $X\in\mod\Lambda$ with $\dim_kX\le n'$, the projective dimension $\projdim_AX$ is either $\infty$ or at most $m$.
\end{proposition} 

Now we are ready to prove Proposition \ref{complexity one}.

\begin{proof}[Proof of Proposition \ref{complexity one}]
Assume that each $B$-module has complexity at most one.

(1) First, we prove $\gldim A$ is finite.
Otherwise, there exists a simple $A$-module $S$ such that $\projdim_AS=\infty$. Take a minimal projective resolution
\[\cdots\to P_2\xrightarrow{f_2} P_1\xrightarrow{f_1} P_0\to S\to0\]
of $S$ in $\mod^{\Z}B$. Take $N>0$ such that $\dim_k\Omega_B^i(S)\le N$ for all $i\ge0$.

Since $\Omega_B^i(S)_0=\Omega_A^i(S)\neq0$ and $\Omega_B^i(S)$ is indecomposable, the inequality $\dim_k\Omega_B^i(S)\le N$ implies $\Omega_B^i(B)_N=0$. Now we take minimal $n\ge1$ such that $\Omega_B^i(S)_{n+1}=0$ for all $i\ge0$. Then each $P_i$ is generated by elements of degrees at most $n-1$, and the sequence
\begin{equation}\label{degree n part}
\cdots\to (P_2)_n\xrightarrow{(f_2)_n} (P_1)_n\xrightarrow{(f_1)_n} (P_0)_n\to S_n=0
\end{equation}
is exact. In particular, each $(P_i)_n\in\inj A$ holds for each $i\ge0$.

We claim that $(f_i)_n$ is a radical map for all $i\ge1$. In fact, let $Q_i$ be a maximal direct summand of $P_i$ such that $Q_i$ is generated in degree $n-1$, and let $g_i:Q_i\to Q_{i-1}$ be the composition $Q_i\subset P_i\xrightarrow{f_i}P_{i-1}\to Q_{i-1}$. Then
\[(f_i)_n=(g_i)_n:(P_i)_n=(Q_i)_n\to (P_{i-1})_n=(Q_{i-1})_n\]
holds. Moreover, the morphism $(g_i)_n$ in $\inj A$ is the image of the morphism $(g_i)_{n-1}$ in $\proj A$ via Nakayama functor. Since $g_i$ is the radical map, so are $(g_i)_{n-1}$ and $(g_i)_n$, as desired.

Let $m$ be the minimal number such that $\Omega_B^m(S)_n\neq0$. Then the exact sequence \eqref{degree n part} shows that $\injdim_A\Omega_B^{m+i}(S)_n=i$ for each $i\ge0$. Since $\dim_k\Omega_B^{m+i}(S)_n\le N$, this contradicts to Proposition \ref{FDC}.

(2) Next, we prove that $A$ is fractionally Calabi-Yau.

We have an exact sequence $0\to DA(-1)\to B\to A\to0$ in $\mod^{\Z}B$. This implies $\Ext^1_{\mod^{\Z/2\Z}B}(A,A)=0$ and hence  $\Ext^1_{\mod^{\Z/2\Z}B}(\Omega_B^i(A),\Omega_B^i(A))=0$ for each $i\ge0$. 
Since $\dim_k\Omega_B^i(A)$ is bounded, Proposition \ref{voigt} shows that there exists $n\ge1$ such that $\Omega_B^n(A)\simeq A$ in $\mod^{\Z/2\Z}B$. Replacing $n$ by its multiple, we can assume that $\Omega_B^n(P)\simeq P$ in $\mod^{\Z/2\Z}B$ for each indecomposable projective $A$-module $P$.
Then there exists $\ell_P\in 2\Z_{>0}$ such that $\Omega_B^n(P)\simeq P(-\ell_P)$ in $\mod^{\Z}B$.
Since $A$ is ring indecomposable, $\ell_P$ is independent of $P$. Thus $\Omega_B^n(A)\simeq A(-\ell)$ holds in $\mod^{\Z}B$. Using the commutative diagram \eqref{serre}, we have $A[-n]\simeq\nu^{-\ell}(A)[-\ell]$ in $\DDD^
{\bo}(\mod A)$. Thus $A$ is twisted fractionally Calabi-Yau.
\end{proof}
}

%%%%%%%%%%%%%%%%%%%%%%%%%%%%%%%%%%%%%%%%%%%%%%%%%%%%%%%%%%%%%%%%%%%
%%%%%%%%%%%%%%%%%%%%%%%%%%%%%%%%%%%%%%%%%%%%%%%%%%%%%%%%%%%%%%%%%%%
%%%%%%%%%%%%%%%%%%%%%%%%%%%%%%%%%%%%%%%%%%%%%%%%%%%%%%%%%%%%%%%%%%%
%%%%%%%%%%%%%%%%%%%%%%%%%%%%%%%%%%%%%%%%%%%%%%%%%%%%%%%%%%%%%%%%%%%

\subsection{\new{Twisted periodicity and $d$-representation finiteness}}

%\subsection{Canonical construction}

An essential part of \cite[2.9]{DI} can be stated as follows.

\begin{proposition}\label{DI2.9}
  Let $\Lambda$ be a twisted $\frac{m}{\ell}$-Calabi--Yau algebra with $d\ge \gl\Lambda$
  for positive integers $\ell,m,d$.  Define $g={\rm gcd}(\ell+m,d+1)$ and $s=\frac{(d+1)\ell-(\ell+m)}{g}=\frac{d\ell-m}{g}$. 
Then $\stmod^\Z T(\Lambda)$ has an \emph{$(s)$-stable} $d$-cluster tilting subcategory $\add\{\Omega_{T(\Lambda)}^{(d+1)i}(\Lambda(i))\mid i\in\Z\}$.
In particular, the following hold.
\begin{enumerate}[(1)]
\item $T(\Lambda)\oplus \bigoplus_{i=1}^s \Omega_{T(\Lambda)}^{(d+1)i}(\Lambda(i))$ is a $d$-cluster tilting object in $\mod^{\Z/s\Z} T(\Lambda)$.

\item If $d+1=(m+\ell)r$ for some $r$, then $T_{r\ell-1}(\Lambda)\oplus \Lambda$ is a $d$-cluster tilting $T_{r\ell-1}(\Lambda)$-module. \marginpar{\color{black}$T_s(\Lambda)=\widehat{\Lambda}/\nu_{\widehat{\Lambda}}^s$}
\end{enumerate} 
 
\end{proposition}
\begin{proof}
Under the equivalence \eqref{happel}, the subcategory in the statement corresponds to a $d$-cluster tilting subcategory $\UU_d(\Lambda)=\add\{\nu_d^i(\Lambda)\mid i\in\Z\}$ of $\DDD^{\bo}(\mod\Lambda)$.
We need to show that it is closed under $(s)$. Since
\[(s)\stackrel{\eqref{serre}}{=}\nu^s[s]=(\nu[-d])^{-\frac{m+\ell}{g}}(\nu^\ell[-m])^{\frac{d+1}{g}}\]
and $\UU_d(\Lambda)$ is closed under both $\nu_d$ and $\phi^*=\nu^\ell[-m]$, the first assertion follows.

(1): This follows by \cite[2.14(a)]{DI} and that $\bigoplus_{i=1}^s (T(\Lambda)\oplus \Omega_{T(\Lambda)}^{(d+1)i}(\Lambda(i))$ is the additive generator of the subcategory obtained by applying the push-down functor $\mod^{\Z}T(\Lambda)\to  \mod^{\Z/s\Z}T(\Lambda)$ on the $(s)$-stable $d$-cluster tilting subcategory $\add\{\Omega_{T(\Lambda)}^{(d+1)i}(\Lambda(i))\mid i\in\Z\}$.

(2):  First note that $\gcd(d+1,m+\ell)=m+\ell$ and $s=r\ell-1$ under the assumption.  Hence, we have
\[
\Omega_{T(\Lambda)}^{(d+1)i}(\Lambda(i)) = \Omega_{T(\Lambda)}^{(m+\ell)ri}(\Lambda)(i) = \Lambda((1-r\ell)i) = \Lambda(-si),
\]
which is $\Lambda$ in $\mod^{\Z/s\Z} T(\Lambda)$.

Since there is an equivalence $\mod^{\Z/s\Z}T(\Lambda)\simeq \mod T_s(\Lambda)$, the claim follows from (1) as $\bigoplus_{i=1}^s T(\Lambda)(i)$ is sent to $T_s(\Lambda)$.
\end{proof}



\iffalse
\begin{proposition}\label{weaker-boolean}
  $\mod^{\Z}B$ has a $(2)$-stable $(n+2)$-cluster subcategory
  \[\CC=\add\{\Omega_{B}^{(n+3)i}(A(i)),\ B(i)\mid i\in\Z\}
  = \{A(2i),\ B(i)\mid i\in\Z\} .\]
In particular, $\mod^{\Z/2\Z} B$ has a $(n+2)$-cluster tilting module $B\oplus B(1)\oplus A$, and $T_2(A)$ has a $(n+2)$-cluster tilting module $T_2(C\oplus A)$.
\end{proposition}

\begin{proof}
For $\Lambda=A$, the assumptions in Proposition \ref{DI2.9} are satisfied by $a=3$, $b=n$ and $d=n+2$. Since $g=n+3$ and $m=2$, the assertion follows.

For the last part, it follows from Lemma \ref{lem:periodicity} that $\Omega_B^{(n+3)i}(A(i))=A(-2i)$. Since $B\oplus B(1)\oplus A$ is the image of $\CC$ via the forgetful functor $\mod^{\Z}B\to\mod^{\Z/2\Z}B$, it is an $(n+2)$-cluster tilting object by \cite[2.14(a)]{DI}.
%sending $\mathcal{C}\subset \mod^{\Z} B$ to $\mod^{\Z} B/(2) \simeq \mod^{\Z/2\Z} B$.
Since there is an equivalence $\mod^{\Z/2\Z}B\simeq\mod T_2(A)$ sending $B\oplus B(1)$ to $T_2(A)$, the last assertion follows.
\end{proof}
\fi

\new{
Now we are ready to prove Theorem \ref{main theorem 2 again}.

\begin{proof}[Proof of Theorem \ref{main theorem 2 again}]
(i)$\Leftrightarrow$(ii) was shown in Proposition \ref{thm:hani-gr}, (ii)$\Rightarrow$(iii) is obvious, (iii)$\Rightarrow$(v) was shown in Proposition \ref{complexity one}, and (v)$\Rightarrow$(ii) was shown in Proposition \ref{tw-fCY to tw-pe}. 

On the other hand, (v)$\Rightarrow$(iv) follows from Proposition \ref{DI2.9}, and (iv)$\Rightarrow$(iii) was shown in \cite{EH}.
Thus all the conditions are equivalent.
\end{proof}
}


%%%%%%%%%%%%%%%%%%%%%%%%%%%%%%%%%%%%%%%%%%%%%%%%%%%%%%%%%%%%%%%%%%%
%%%%%%%%%%%%%%%%%%%%%%%%%%%%%%%%%%%%%%%%%%%%%%%%%%%%%%%%%%%%%%%%%%%
%%%%%%%%%%%%%%%%%%%%%%%%%%%%%%%%%%%%%%%%%%%%%%%%%%%%%%%%%%%%%%%%%%%
%%%%%%%%%%%%%%%%%%%%%%%%%%%%%%%%%%%%%%%%%%%%%%%%%%%%%%%%%%%%%%%%%%%


\section{\new{When is $T(A)$ $d$-representation finite?}}

\new{
Let $\TT$ be a $k$-linear Hom-finite Krull-Schmidt triangulated category with Serre functor $\nu$. 
Then there exists a bijection between the isomorphism classes of basic silting objects in $\TT$ and the bounded-co-t-structures of $\TT$ \cite{KY}. For a silting object $P$, the corresponding co-t-structure is given by
\[\TT_{\ge0}=\cdots*(\add P[-1])*(\add P)\ \mbox{ and }\ \TT_{\le0}=(\add P)*(\add P[1])*\cdots.\]
We fix a positive integer $d$. 

\begin{definition}
We call a silting object $P\in\TT$ \emph{$d$-silting} if $\nu_d(\TT_{\ge0})\subset\TT_{\ge0}$ holds, or equivalently, $\nu_d^{-1}(\TT_{\le0})\subset\TT_{\le0}$ holds.
\end{definition}

\begin{corollary}\label{from silting to cluster tilting}
Let $\Lambda$ be a finite-dimensional $k$-algebra of finite global dimension, and $P$ a silting complex.
If $\Lambda$ is $\nu_d$-finite and $P$ is $d$-silting, then $\DDD^{\bo}(\mod\Lambda)$ has a locally bounded $d$-cluster tilting subcategory
\[\UU_d(P):=\add\{\nu_d^i(P)\mid i\in\Z\}.\]
\end{corollary}
}



\begin{lemma} \label{tiltandsilt}
  Let $r$ be a positive integer, $\Lambda$ and $\Lambda'$ finite-dimensional $k$-algebras of finite global dimension. Assume that $X$ is a $\Lambda$-module such that
  \begin{enumerate}\renewcommand{\theenumi}{\roman{enumi}}
  \item  $\bigoplus_{i=0}^{r-1}\nu^i(X)\in\KKK^{\bo}(\proj\Lambda)$ is a tilting object;
  \item $\Hom_{\KKK^{\bo}(\proj\Lambda)}(\nu^i(X),X) = 0$ for all $i\in\{1,\ldots, r-1\}$. 
  \end{enumerate}
  Set $Q = X\otimes\Lambda'\in\mod\left(\Lambda\otimes_k\Lambda'\right)$. 
  Then $\bigoplus_{i=0}^{r-1}\nu^i(Q)$ %\in\KKK^{\bo}(\proj\Lambda\otimes\Lambda')$
  is a tilting object, and
$\bigoplus_{i=0}^{r-1}\nu^i(Q)[i]$ %\in\KKK^{\bo}(\proj\Lambda\otimes\Lambda')$
  a silting object, in $\KKK^{\bo}(\proj\left(\Lambda\otimes_k\Lambda'\right))$.
\end{lemma}
%
%
%% \begin{lemma} \label{tiltandsilt}
%%   Let $\Lambda$ and $\Lambda'$ be finite-dimensional $k$-algebras of finite global dimension, and $X$ a $\Lambda$-module satisfying the following conditions:
%%   \begin{enumerate}\renewcommand{\theenumi}{\roman{enumi}}
%%   \item $2|X|=\Lambda$;
%%   \item $\Ext^l(X,X) = 0$ for all $l\ge 1$;
%%   \item $\Hom_{\KKK^{\bo}(\proj\Lambda)}(\nu(X),X[l]) = 0$ for all $l\in\{0,1,\ldots,2\gl\Lambda\}$.
%%   \end{enumerate}
%%   Set $Q = X\otimes\Lambda'\in\mod\left(\Lambda\otimes_k\Lambda'\right)$. 
%%   Then $Q\oplus\nu(Q)$ is a tilting object, and $Q\oplus\nu(Q)[1]$ a silting object, in $\KKK^{\bo}(\proj\left(\Lambda\otimes_k\Lambda'\right))$.
%% \end{lemma}

\begin{remark}
In particular, Lemma~\ref{tiltandsilt} applies in the following situation:
$\Lambda$ is a hereditary, and $X\in\mod\Lambda$ is such that $r|X| = |\Lambda|$, $\Ext^1(X,X) = 0$, and $\Hom_{\KKK^{\bo}(\proj\Lambda)}(\nu^i(X),X[\ell]) = 0$ holds for all $1\le i \le r-1$ and $0\le \ell\le i+1$.%
\marginpar{Enough $\lfloor i/2\rfloor\le \ell\le i+1$?}
\end{remark}

\new{
\begin{proof}
  Set $M= \bigoplus_{i=0}^{r-1}\nu^i(X)$.
  We need to establish the following:
\begin{enumerate}[a)]
\item $\Hom_{\kbpl}(X,X[\ell]) = 0$ for all $\ell\in\Z\setminus\{0\}$;
\item $\Hom_{\kbpl}(\nu^i(X),X[\ell])=0$ for all $i\in\{1,\ldots,r-1)$, $\ell\in\Z$;
\item $\Hom_{\kbpl}(X,\nu^j(X)[\ell])=0$ for all $j\in\{1,\ldots,r-1)$, $\ell\in\Z\setminus\{0\}$;
\item $\thick M  = \kbpl$.
\end{enumerate}

(a)
This is immediate from the assumptions that $\Lambda$ is hereditary and $\Ext^1(X,X)=0$.

(b)
The identity $\Hom_{\kbpl}(\nu^i(X),X[\ell])=0$ holds for all $\ell<0$, since $\nu^i(X)\in\DDD^{\le0}$ and, by assumption, also for $0\le\ell \le i+1$. 
Let $\ell\ge i+2$. Then
\[\Hom_{\kbpl}(\nu^i(X),X[\ell]) \simeq \Hom(\tau^i(X), X[\ell-i]) =0 \]
since $\tau^i\in\DDD^{\ge0}$, $X[\ell-i]\in\DDD^{\le2}$ and $\Lambda$ is hereditary.

(c)
We have
\[ \Hom_{\kbpl}(X,\nu^j(X)[\ell]) \simeq D \Hom_{\kbpl}(\nu^{j-1}(X), X[-\ell]) =0 ,\]
where the last equality follows from (b) if $j>1$ and from (a) if $j=1$.

(d)
Inserting $\ell=0$ into statement (b) gives $\Hom_{\kbpl}(\nu^i(X),X)=0$ for all $i\in\{1,\ldots,r-1\}$, so $\nu^i(X)$ and $\nu^j(X)$ have no common direct summands for $i\ne j$.
Thus
$|M|=\sum_{i=0}^{r-1}\left|\nu^i(X)\right| = r|X| = |\Lambda|$.
Moreover, from (a)--(c) it follows that
\[\Hom_{\kbpl}\left(M, M[\ell]\right) = 0,\quad \mbox{for all}\; \ell\in\Z.\]
Since $\Lambda$ is hereditary, this implies \cite{} that $\thick M = \kbpl$.
\end{proof}
}

%% In particular, the conditions (ii)--(iii) in Lemma~\ref{tiltandsilt} are satisfied if $X$ is a projective module such that % $2|X|=|\Lambda|$ and
%% $\Ext^l(\nu(X),X)=0$ for all $l\in\{0,1,\ldots,\gl\Lambda\}$. 

\begin{proof}[Proof of Lemma~\ref{tiltandsilt}]
  \ldots
\end{proof}

\old{
\begin{proof}
  First note that the condition (iii) ($l=0$) implies that $Q$ and $\nu(Q)$ have no non-trivial direct summands in common. Hence,
  \[|Q\oplus\nu(Q)| = 2|Q| = 2|X||\Lambda'| = |\Lambda||\Lambda'| = |\Lambda\otimes_k\Lambda'| \,.\]
  To show that $Q\oplus\nu(Q)$ is a tilting complex, it remains to prove that
  $\Hom_{\KKK^{\bo}(\proj\left(\Lambda\otimes_k\Lambda'\right))} (Q\oplus\nu(Q),(Q\oplus\nu(Q))[\ne0]) = 0$.

The restriction functor
\[\res : \KKK^{\bo}(\proj\left(\Lambda\otimes_k\Lambda'\right)) \to \KKK^{\bo}(\proj\Lambda) \]
is a faithful triangle functor, with left adjoint
%
\[-\otimes_k\Lambda': \KKK^{\bo}(\proj\Lambda) \to \KKK^{\bo}(\proj(\Lambda\otimes_k\Lambda'))\, . \]
%
From the conditions (ii)--(iii), it is straightforward to verify that $\Hom_{\KKK^{\bo}(\proj\Lambda)}(X\oplus\nu(X),(X\oplus\nu(X))[\ne0])=0$. 
% is a tilting object in $\KKK^{\bo}(\proj\Lambda)$. 
Now, for all $l\in\Z\setminus\{0\}$,
\begin{align*}
  &\Hom_{\KKK^{\bo}(\proj\left(\Lambda\otimes_k\Lambda'\right))}(\nu(Q),\nu(Q)[l]) \simeq
  \Hom(Q,Q[l]) \simeq
  \Hom(X\otimes\Lambda',(X\otimes\Lambda')[l]) \\
  &\simeq \Hom_{\KKK^{\bo}(\proj\Lambda)} (X, \res(X\otimes\Lambda')[l]) \simeq \Hom_{\KKK^{\bo}(\proj\Lambda)} \left(X, X^{\oplus\dim_k\Lambda'}[l]\right) = 0,\\
  \intertext{and}
  &\Hom_{\KKK^{\bo}(\proj\left(\Lambda\otimes_k\Lambda'\right))}(Q,\nu(Q)[l]) \simeq
  D\Hom(Q,Q[-l]) = 0.
\end{align*}
%
Moreover, for all $l\in\Z$, we have
\begin{multline}\label{homzero}
  \Hom_{\KKK^{\bo}(\proj\left(\Lambda\otimes_k\Lambda'\right))}(\nu(Q),Q[l]) \simeq
  \Hom(\nu(X)\otimes D\Lambda', (X\otimes \Lambda')[l]) \\
  \subset
  \Hom_{\KKK^{\bo}(\proj\Lambda)} (\res(\nu(X)\otimes D\Lambda') , \res(X\otimes\Lambda')[l]) \simeq
  \Hom(\nu(X)^{\oplus\dim_kD\Lambda'}, X^{\oplus\dim_k\Lambda'}[l]) = 0 .
 %%  \Hom(Q,\nu^{-1}(Q)[l])  
 %%  \simeq \Hom_{\KKK^{\bo}(\proj\Lambda)} (X,\res(\nu^{-1}(Q))[l]) \\ \simeq
 %%  \Hom(\nu(X),\res(Q)[l]) \simeq \Hom\left(\nu(X),Q^{\oplus\dim_k\Lambda'}[l]\right)
 %% =0 .
\end{multline}
%
Hence, \[\Hom_{\KKK^{\bo}(\proj(\Lambda\otimes\Lambda'))}(Q\oplus\nu(Q),(Q\oplus\nu(Q))[\ne0])=0,\]
so $Q\oplus\nu(Q)$ is a tilting complex. 

We proceed to show that $Q\oplus\nu(Q)[1]$ is silting.
Notice first that $Q\oplus\nu(Q)\in\thick(Q\oplus\nu(Q)[1])$ and thus, since $Q\oplus\nu(Q)$ is tilting,
$\thick(Q\oplus\nu(Q)[1]) = \KKK^{\bo}(\proj\left(\Lambda\otimes_k\Lambda'\right))$. 
Let $l$ be a positive integer. The identities
\[\Hom(Q,(Q\oplus\nu(Q)[1])[l]) = \Hom(\nu(Q)[1],(\nu(Q)[1])[l]) = 0 \]
follow from $Q\oplus\nu(Q)$ being tilting, and
\[\Hom(\nu(Q)[1],Q[l])\simeq \Hom(\nu(Q),Q[l-1]) = 0\]
from Equation~\eqref{homzero}.
%
\end{proof}
}

\begin{lemma} \label{CYcot}
  Let $A$ be $\frac{d-r}{r+1}$-Calabi--Yau, where $d>r\ge1$, and let $Q\in\KKK^{\bo}(\proj A)$ be such that $M= \bigoplus_{i=0}^{r-1} \nu^i(Q)$ is a tilting object and
  $N = \bigoplus_{i=0}^{r-1} \nu^i(Q)[i]$ is a silting object.
  Then $N$ is $d$-silting.
%  $\nu_{d}(\TT_{\ge0})\subset \TT_{\ge0}$ holds. 
\end{lemma}

\begin{proof}
It suffices to show that $\Hom(\nu_d(N),N[>0])=0$.
First note that
\[\mbox{$A$ is $\frac{d-r}{r+1}$-Calabi--Yau} \;\Leftrightarrow\; [d-r]\simeq \nu^{r+1} \;\Leftrightarrow\; \nu_d \simeq \nu^{-r}\circ[-r] \,. \]
For ease of notation, set $F = \nu\circ[1]$. 
Let $l$ be a positive integer, and $0\le i<r$. Using the above isomorphism of functors, and the assumption that $M$ is tilting, we get
\begin{align*}
  &\Hom(\nu_d(F^i(Q)),F^i(Q)[l]) \simeq \Hom(\nu_d(Q),Q[l]) 
  \simeq \Hom(\nu^{-r}(Q)[-r],Q[l]) \\
  &\simeq \Hom(Q,\nu^r(Q)[l+r]) \simeq D\Hom(\nu^{r-1}(Q)[l+r], Q) = 0. \\
\intertext{For $i<j<r$:}
  &\Hom(\nu_d(F^i(Q)),F^j(Q)[l]) \simeq \Hom(\nu(Q)[-d],F^{j-i}Q[l])\simeq \Hom(\nu(Q),\nu^{j-i}Q[(j-i)+l+d]) \\
  &\simeq \Hom(Q,\nu^{j-i-1}Q[(j-i)+l+d]) = 0, \\
  \intertext{as $\nu^{j-i-1}(Q)\in\add M$ and $j-i+l+d>0$. Moreover,}
  &\Hom(\nu_d(F^j(Q)), F^i(Q)[l]) \simeq \Hom(\nu^{-r}(Q)[-r], F^{i-j}Q[l]) \simeq \Hom(Q, \nu^{r+i-j}(Q)[l+r+i-j]) =0,
\end{align*}
since $\nu^{r+i-j}(Q)\in \add M$ and $l+r+i-j>0$. This concludes the proof.
\end{proof}


Combining Lemma~\ref{tiltandsilt} and Lemma~\ref{CYcot}, we get the following result.

\begin{proposition} \label{CYsiltCT}
  Let $\Lambda$ and $\Lambda'$ be finite-dimensional $k$-algebras of finite global dimension, such that the algebra $A=\Lambda\otimes_k\Lambda'$ is $\frac{b}{a}$-Calabi--Yau, with $a\ge2$.
  If $X$ is a $\Lambda$-module satisfying the conditions (i)--(ii) in Lemma~\ref{tiltandsilt} with $r = a-1$ and $d = a+b-1$, then
  $X\otimes\Lambda' \,\oplus\, T(A)$ is an $(a+b-1)$-cluster-tilting module of $T(A)$.
\end{proposition}

Note that, under the conditions of Proposition~\ref{CYsiltCT}, the functorial isomorphism $\Omega^{a+b}\simeq (-a)$ holds in $\stmod^{\Z}T(A)$, by Proposition~\ref{periodic to fCY}.

%% \begin{proposition} \label{CYsiltCT}
%%   Let $\Lambda$ and $\Lambda'$ be finite-dimensional $k$-algebras of finite global dimension, such that the algebra $A=\Lambda\otimes_k\Lambda'$ is $\frac{d-2}{3}$-Calabi--Yau.
%%   If $X$ is a $\Lambda$-module satisfying the conditions (i)--(iii) in Lemma~\ref{tiltandsilt}, then
%%   $X\otimes\Lambda' \,\oplus\, T(A)$ is a $d$-cluster-tilting module of $T(A)$, and $\Omega^{d+1}\simeq (-3)$ holds in $\stmod^{\Z}T(A)$.
%% \end{proposition}

\begin{proof}
For $r = a-1$ and $d = a+b-1$, set
\[Q = X\otimes \Lambda',\quad
M = \bigoplus_{i=0}^{r-1}\nu^i(Q), \quad\mbox{and}\quad
N = \bigoplus_{i=0}^{r-1}\nu^i(Q)[i].\]
By Lemma~\ref{tiltandsilt}, $N$ is a silting object in $\KKK^{\bo}(\proj A)\simeq\stmod^{\Z}T(A)$, and Lemma~\ref{CYcot} implies that the co-$t$-structure $(\TT_{\ge0},\TT_{\le0})$ in $\KKK^{\bo}(\proj A)$ given by this silting object satisfies $\nu_d(\TT_{\ge0})\subset \TT_{\ge0}$.
%
By the Calabi--Yau property,
\[ \nu_d^{a}\simeq \nu^a\circ[-da]\simeq [b-(a+b-1)a] = [-(a-1)(a+b)] \]
and, since $\gl A<\infty$, this implies that $A$ is $\nu_d$-finite. 
Corollary~\ref{from silting to cluster tilting} now gives that $\UU:=\UU_d(N)\subset\KKK^{\bo}(\proj A)$ is a locally bounded $d$-cluster-tilting subcategory.

Again by the Calabi--Yau property,
\[(\nu\circ[1])^r \simeq \nu^{a-1}\circ[a-1] \simeq \nu^{-1}\circ[b]\circ[a-1] = \nu^{-1}\circ[d] \simeq \nu_d^{-1},\]
and hence
\[\textstyle \UU = \add\left\{\nu_d^i\left(\bigoplus_{i=0}^{r-1}\nu^i(Q)[i]\right) \mid i\in\Z\right\} = \add\{(\nu\circ[1])^j(Q) \mid j\in\Z \} \,.\]
In particular, $\UU$ is invariant under the auto-equivalence $\nu\circ[1]$ of $\KKK^{\bo}(\proj A)$ or, equivalently, under degree shift $(1)$ in the category $\stmod^{\Z} T(A)$.
A fundamental domain for the action of the induced permutation of $\ind\UU$ is given by $\ind(\add Q)$.
From \cite[Theorem~2.3]{DI} it now follows that $Q\oplus T(A)\in\mod T(A)$ is a $d$-cluster-tilting module.

%% Let $\mathcal{V}\subset\mod^{\Z}T(A)$ be the preimage of $\UU_d(Q\oplus\nu(Q)[1])$ under the natural functor $\mod^{\Z}T(A) \to \stmod^{\Z}T(A)$. By \cite[Lemma~3.1(b)]{DI}, this is a 
\end{proof}

\begin{corollary}
  Let $A = (k\mathbb{A}_2)^{\otimes l} \otimes_k (k\mathbb{D}_4)^{\otimes m}$ for some $l,m\ge0$, and set $d=l+2m+2$. Then $T(A)$ is $d$-representation-finite, and $\Omega^{d+1} \simeq (-3)$ in $\stmod^\Z T(A)$.
\end{corollary}

\begin{proof}
  As $k\mathbb{A}_2$ and $k\mathbb{D}_4$ are $\frac{1}{3}$ respectively $\frac{2}{3}$ fractionally Calabi--Yau, it follows that $A$ is $\frac{l+2m}{3}$-Calabi--Yau \cite[Proposition~1.4]{HerIya}.
  \marginpar{Assumption on $k$?}
To apply Proposition~\ref{CYsiltCT}, we need to show that each of $k\mathbb{A}_2$ (for the case $l\ge1$) and $k\mathbb{D}_4$ (for $l=0$) has a module $X$ satisfying the conditions (i) and (ii) in Lemma~\ref{tiltandsilt}.

For the algebra $k\mathbb{A}_2$, let $X$ be an indecomposable projective module. 
%
In the case of $\mathbb{D}_4$, let let $X = P_1\oplus P_2\in \mod kQ$, where $Q$ is the following quiver:
\[ \xymatrix{ &1 & &\\ 2\ar[r] & 3\ar[u] & \ar[l]4} \]
%
In both cases, it is straightforward to verify that $X$ satisfies the conditions
in Lemma~\ref{tiltandsilt}.
\end{proof}


\begin{corollary}
  Let $A = (k\mathbb{A}_m) \otimes_k \Lambda'$, where $\Lambda'$ is $\frac{b}{m+1}$-Calabi--Yau, and set $d = 2m+b-1$. 
  Then $T(A)$ is $d$-representation-finite, and $\Omega^{d+1} \simeq (-(m+1))$ in $\stmod^\Z T(A)$.
\end{corollary}

\begin{proof}
  We may assume the quiver $\mathbb{A}_m$ to have linear orientation. Let $X\in\mod k\mathbb{A}_m$ be the simple projective module.
  Then $M = \bigoplus_{i=0}^{m-1} \nu^i(X)$ is a tilting object \texttt{[reference?]}.
  For any $i\ge1$, $\nu^{-i}(X)\in\DDD^{>0}$ and thus
  \[ \Hom(\nu^i(X),X)\simeq \Hom(X, \nu^{-i}(X))\simeq 0 \,.\]
  Hence $X$ satisfies the conditions (i) and (ii) in Lemma~\ref{tiltandsilt} for $r=m$.

  With $k\mathbb{A}_m$ and $\Lambda'$ being fractionally Calabi--Yau of dimensions $\frac{m-1}{m+1}$ respectively $\frac{b}{m+1}$, the algebra $A = k\mathbb{A}_m\otimes_k \Lambda'$ is $\frac{m+b-1}{m+1}$-Calabi--Yau.
  Applying Proposition~\ref{CYsiltCT}, we get that $Q=X\otimes\Lambda'\,\oplus\, T(A) \in\mod T(A)$ is a $d$-cluster-tilting module, for $d = (m+1) + (m+b-1) -1 = 2m+b-1$.
  By Proposition~\ref{periodic to fCY}, the functorial isomorphism $\Omega^{2m+b}\simeq (-1-m)$ holds is $\stmod T(A)$. 
\end{proof}

\begin{example}
  The algebra $\Lambda'= (k\mathbb{A}_m)^{\otimes(n-1)}$ is $\frac{(n-1)(m-1)}{m+1}$-Calabi--Yau, so $T\left((k\mathbb{A}_m)^{\otimes(n-1)}\right)$ is $d$-representation-finite for $d=(n+1)(m-1) +1$, and $\Omega^{(n+1)(m-1)+2}\simeq (-(m+1))$ in $\stmod T\left((k\mathbb{A}_m)^{\otimes(n-1)}\right)$. 
\end{example}


%% Since $A$ is $\frac{n}{3}$-Calabi--Yau, it is $\nu_{n+2}$-finite.
%% Thanks to Corollary \ref{from silting to cluster tilting}, it suffices to find out a silting complex $P$ of $A$ satisfying $\nu_{n+2}(\TT_{\ge0})\subset\TT_{\ge0}$.

%% Let $e_0$ and $e_1$ be primitive idempotents of $k\mathbb{A}_2$. Regard $A=k\mathbb{A}_2\otimes_kA'$, and
%% let $Q=A(e_1\otimes 1_{A'})\in\proj A$.

%% \begin{proposition}
%% $P:=Q\oplus \nu(Q)[1]$ a silting complex of $A$.
%% \end{proposition}
%% \begin{proof}
%% To see that $Q\oplus \nu(Q)[1]$ is silting, perform the follow sequence of mutations.

%% For $i\in \{0, 1,\ldots, n-1\}$, let 
%% \[
%% M_i:=\bigoplus_{\substack{\pmb{x}\,:\, |\pmb{x}|=i\\x_1=0}} P_{\pmb{x}}.
%% \]
%% Then $M_0\oplus M_1\oplus \cdots \oplus M_{n-1} = A/Q$.

%% Similarly, for $i\in\{1,2,\ldots, n\}$, consider the direct sum
%% \[
%% N_i := \bigoplus_{\substack{\pmb{x}\,:\, |\pmb{x}|=i\\x_1=1}} I_{\pmb{x}}
%% \]
%% of injective $A$-modules.
%% Note that $N_1\oplus N_2\oplus \cdots N_{n} = \nu(Q)$.

%% \begin{itemize}
%% \item[(0)] Mutate at $A/Q$ yields $A/Q[1]\oplus Q$ (since $\Hom_A(A/Q,Q)=0$).

%% \item[(1)] Mutate at $M_{n-1}[1]$ for $n-1$ times yields $((A/Q)/M_1)[1] \oplus N_1[1] \oplus Q$.
%% \\ $\vdots$

%% \item[($i$)] Mutate at $M_{n-i}[1]$ for $n-i$ times yields $((A/Q)/(M_{n-1}\oplus\cdots M_{n-i}))[1] \oplus (N_1 \oplus \cdots \oplus N_{i})[1] \oplus Q$
%% \\ $\vdots$

%% \item[($n-1$)] Mutate at $M_1[1]$ yields $(N_1 \oplus \cdots \oplus N_{n-1} \oplus M_0)[1]\oplus Q$.  Note that $M_0=N_n$, so we have $\nu(Q)[1]\oplus Q$ now.
%% \end{itemize}
%% \end{proof}

%% \begin{question}
%% Does the condition $\nu_d(\TT_{\ge0})\subset\TT_{\ge0}$ hold?
%% This is equivalent to $\Hom_{\TT}(\nu_d(P),P[>0])=0$, since $(\TT_{\ge0},\TT_{\le0})$ is a t-structure.
%% \end{question}


%%%%%%%%%%%%%%%%%%%%%%%%%%%%%%%%%%%%%%%%%%%%%%%%%%%%%%%%%%%%%%%%%%%
%%%%%%%%%%%%%%%%%%%%%%%%%%%%%%%%%%%%%%%%%%%%%%%%%%%%%%%%%%%%%%%%%%%
%%%%%%%%%%%%%%%%%%%%%%%%%%%%%%%%%%%%%%%%%%%%%%%%%%%%%%%%%%%%%%%%%%%
%%%%%%%%%%%%%%%%%%%%%%%%%%%%%%%%%%%%%%%%%%%%%%%%%%%%%%%%%%%%%%%%%%%
\section{Detour: Serre-formal modules and periodicity}

\marginpar{rewrite or remove this section?}

Before showing the missing arrow in \eqref{dgm:implications}, let us explain how one can find out the terms of the projective resolution of the simple $B$-modules in a certain special case.  In particular, this lets us understand the terms appearing in the projective resolution of $B$ as a $B$-$B$-bimodule, which inspired our argument in Section \ref{sec:fCY to per}.  We will work in a slightly greater generality.

Suppose $\Lambda$ is a Iwanaga-Gorenstein algebra, i.e. $\injdim {}_{\Lambda}\Lambda$ and $\injdim \Lambda_\Lambda$ are both finite.
As in subsection \ref{subsec:CYdim}, we denote by $\nu:=-\Lotimes_\Lambda D\Lambda$ the Serre functor on $\Kb(\proj\Lambda)$.

\begin{definition}\label{def:nuformal}
An indecomposable $\Lambda$-module $M$ of is called \emph{positively Serre-formal} if
\begin{enumerate}
\item it has finite projective dimension $\projdim X<\infty$, and
\item for every positive integer $m>0$, there is some $\ell\in \Z$ such that the cohomology of $\nu^m(M)$ is concentrated in a single degree.
\end{enumerate}
\emph{Negatively Serre-formal} modules are defined dually via finite injective dimension and power of inverse Serre functor.  
Denote by $M^{\prec m}$ and $M^{\succ m}$ the (stalk) cohomology of $\nu^m(M)$ and $\nu^{-m}(M)$ for all $m\geq 0$ respectively.  For a positively (resp. negatively) Serre-formal module $M$, if $M^{\prec m}\cong M$ (resp. $M^{\succ m}$) for some positive integer $m>0$, then we say that $M$ is \emph{periodically Serre-formal}.

A Iwanaga-Gorenstein algebra ${}_\Lambda\Lambda$ is \emph{positively (resp. negatively, resp. periodically) Serre-formal} if so are all of the indecomposable direct summands of the regular representation ${}_\Lambda\Lambda$.
\end{definition}

For simplicity, we may write `$\nu^+$-formal' and `$\nu^-$-formal' instead of positively and negatively Serre-formal respectively.  Algebras or modules that are both $\nu^+$-formal and $\nu^-$-formal are said to be \emph{Serre-formal} or simply $\nu^\pm$-formal.

Serre-formality (for algebras) was first introduced in the spiritual prequel \cite{CIM}.  The toy model for Serre-formal algebras and modules come from hereditary algebras and its higher generalisation - see \cite{CIM} and \cite{HIO} for the relevant terminologies.
\begin{example}
Let $d$ be a positive integer and assume $A$ is a $d$-hereditary algebra.
Then for any $i\in \Z$, all $d$-preprojective, $d$-preinjective, $d$-regular modules are $\nu^\pm$-formal.  Moreover, $d$-preprojective and $d$-preinjective modules are periodically Serre-formal if, and only if, $A$ is $d$-representation-finite.  On the other hand, $d$-regular modules are always periodically Serre-formal by \cite[Cor 6.23]{HIO}.
\end{example}

A characterisation of periodically Serre-formal algebras is given by the Calabi--Yau property.
\begin{proposition}{\rm \cite[Prop 5.9]{CIM}}\label{nuformal-fCY}
A Serre-formal algebra $A$ is periodically Serre-formal if and only if it is twisted fractionally Calabi--Yau.
\end{proposition}

Unlike $d$-hereditary algebras or $d$-representation-(in)finite algebras, $\nu^\epsilon$-formal algebras (for any $\epsilon\in\{+,-,\pm\}$) are closed under tensor product over the underlying field \cite[Prop 5.4]{CIM}.  Analogously, this holds also for tensor product of Serre-formal modules.  The proof is in verbatim to \cite[Prop 5.4]{CIM} and we leave it as an exercise.
\begin{proposition}
Suppose $M\in \mod\Lambda$ and $M'\in\mod\Lambda'$ are $\nu^\epsilon$-formal with $\epsilon$ being one of the symbols $+, -$, or $\pm$.
Then $M\otimes M'$ is also Serre-formal.
\end{proposition}

The technique used in \cite{CIM} can be ported directly in this setting to write down the projective resolution of $M$ as a graded module over the trivial extension $\Gamma$ of $\Lambda$, which we explain now.

First, recall that for $\Lambda$-modules $M,N$ and a $\Lambda$-linear map $f:D\Lambda\otimes_\Lambda M\to N$, we have a $\Gamma$-module $(M, N, f)$.  A graded $\Gamma$-module is always of this form with $M$ in degree $i$ and $N$ in degree $i+1$, for some $i\in \Z$.  

For a projective $\Lambda$-module $P$, we write $\hat{P}:=\Gamma\otimes_\Lambda P = (P,\nu(P),\mathrm{id}) \in \proj \Gamma$.  For a morphism $(f:P\to P')\in \proj \Lambda$, write $(\hat{f}:\hat{P}\to\hat{P}')\in \proj \Gamma$ given by $\left[\begin{smallmatrix} f & 0 \\ 0& \nu(f)\end{smallmatrix}\right]: P\oplus \nu(P) \to P'\oplus \nu(P')$.

Now we can describe the projective resolution of Serre-formal modules.

\begin{proposition}\label{nuformal}
Suppose $\Lambda$ is an Iwanaga-Gorenstein algebra, $\Gamma:=T(\Lambda)$ is its trivial extension, and $\nu$ is the Serre functor on $\Kb(\proj\Lambda)$.
Let $M\in\mod\Lambda$ be a positively Serre-formal module, then the following holds.
\begin{enumerate}
\item For any $m\geq 0$, we have $\Omega_{\Gamma}^{d_m}(M^{\prec m})\cong M^{\prec m+1}(-1)$ for some integers $d_m>0$.

\item Assume, moreover, that $M$ is periodically Serre-formal and denote by $m$ the minimal positive integer such that $M^{\prec m}\cong M$.  Then there is some $n>0$ such that $\Omega_{\Gamma}^n(M)\cong M(-m)$.  Hence, $M$ is (graded) $\Omega$-periodic as a $\Gamma$-module.
\end{enumerate}
\end{proposition}
\begin{proof}
This is analogous to \cite[Lemma 5.5]{CIM}; we explain below the details for completeness.

Since $\nu$ is an autoequivalence on category of perfect complexes over $\Lambda$, $M^{\succ m}$ is also positively Serre-formal for all $m>0$.
Denote by 
\[
\xymatrix{
 0 \ar[r]& P_{d_m} \ar[r]^{f_{d_m}} & \cdots\ar[r]& P_{d_{m-1}+1} \ar[r]^{f_{d_{m-1}+1}} & P_{d_{m-1}} \ar[r]^{f_{d_{m-1}}} & M^{\succ m}\ar[r] & 0
}
\]
the projective resolution of $M^{\succ m}$ for all $m\geq 0$ with $d_{-1}:=0$.

By Serre-formality of $M^{\prec m}$, applying $\nu$ is exact on such sequence, except at the leftmost term, which has cohomology $M^{\prec m+1}$.  Hence, we obtain exact sequences
\[
\mathcal{P}_{M,m}:=\Big(\xymatrix{
 0 \ar[r]& M^{\succ m+1}(-1)\ar[r]^{\phantom{abcd}\nu(f_{d_m})} & \hat{P}_{d_m} \ar[r]^{\hat{f}_{d_m}} & \cdots \ar[r]& \hat{P}_{d_{m-1}} \ar[r]^{f_{d_{m-1}}} & M^{\succ m}\ar[r] & 0
}\Big)
\]
of graded $\Gamma$-modules for all $m\geq 0$, and (1) follows.

For (2), we glue these exact sequence (with appropriate graded shift) together in the natural way to obtain a projective resolution of the graded $\Gamma$-module $M$, and the claim follows by definition of the periodically Serre-formal.
\end{proof}

Note that for a $\Gamma$-module $M$ over any algebra $\Gamma$, the projective resolution of $M$ as $\Gamma$-module is given by $M\otimes_\Lambda\mathcal{Q}_\Gamma$ where $\mathcal{Q}_\Gamma$ is the projective resolution of $\Gamma$ as $\Gamma^e$-bimodule.

Consider now $A$ is an algebra that has finite global dimension and is periodically Serre-formal algebra.
Then $A$ is $\frac{m}{\ell}$-Calabi--Yau for some integers $m$ and $\ell>0$ by Proposition \ref{nuformal-fCY}.
Suppose every simple module $S$ appears in $\add(\nu^i(A)\mid i\in \Z)$, then the projective resolution $\mathcal{Q}_S$ as a module over the trivial extension $B:=T(A)$ is periodic and is given by the Yoneda composition of $\mathcal{P}_{S,i}$ over $i\in \Z_{\geq 0}$.  

\comment{\texttt{TODO: get terms of $\mathcal{Q}_B$??}}



\iffalse
As a consequence, the minimal Calabi--Yau dimension gives lower bound on the periodicity of $B$.

\begin{corollary}\label{periodicity2}
Let $A$ be a finite-dimensional $\frac{m}{\ell}$-Calabi--Yau algebra of finite global dimension, and $B:=T(A)$ be its trivial extension.
If $\frac{m}{\ell}$ is the minimal Calabi--Yau dimension, then the minimal possible twisted (bimodule) periodicity of $B$ is $m+\ell$.
\end{corollary}
\begin{proof}
Suppose that $\Omega^{s}\simeq (-a)$ on $\stmod^{\Z}B$ for some $s\leq m+\ell$ and some integer $a$.  Then by the same argument as in the proof of Proposition \ref{periodicity}, we have \[
\mathrm{Id}\simeq \Omega^{a+s-a}(-a)\simeq \nu^{a}[s-a],\]
which means that $A$ is $\frac{a-s}{a}$-Calabi--Yau.  By the minimality assumption, we have $a=\ell$ and $s=m+\ell$.
\end{proof}

\begin{lemma}\label{lem:periodicity}
$\Omega^{n+3}\simeq(-3)$ as auto-equivalences on $\stmod^{\Z}B$
\end{lemma}
\begin{proof}
Since our $A$ is $\frac{n}{3}$-Calabi--Yau, the assertion is immediate from Proposition \ref{periodicity}.
\end{proof}
\fi


%%%%%%%%%%%%%%%%%%%%%%%%%%%%%%%%%%%%%%%%%%%%%%%%%%%%%%%%%%%%%%%%%%%
%%%%%%%%%%%%%%%%%%%%%%%%%%%%%%%%%%%%%%%%%%%%%%%%%%%%%%%%%%%%%%%%%%%
%%%%%%%%%%%%%%%%%%%%%%%%%%%%%%%%%%%%%%%%%%%%%%%%%%%%%%%%%%%%%%%%%%%
%%%%%%%%%%%%%%%%%%%%%%%%%%%%%%%%%%%%%%%%%%%%%%%%%%%%%%%%%%%%%%%%%%%

\section{Fractionally Calabi--Yau algebras have periodic trivial extensions}\label{sec:fCY to per}

Throughout this section, let $A$ be a finite dimensional algebra over a field $k$. As usual, let $A^e=A\otimes_kA^{\op}$ be the enveloping algebra. For $i\ge0$, we write
\[(DA)^{\Lotimes i}:=\overbrace{(DA)\Lotimes_A\cdots\Lotimes_A(DA)}^i.\]
Assume that $A$ is $\frac{m}{\ell}$-Calabi--Yau. Then there is an isomorphism
\begin{equation}\label{define s}
s:(DA)^{\Lotimes\ell}\simeq A[m]\ \mbox{ in }\ \DDD(A^e).
\end{equation}

For any $x\in B$, we denote by $(x\cdot)$ the map given by left-multiplying $x$ on $B$.

We consider the composition
\begin{equation}\label{define t}
t:=\Big(DA[m]\xrightarrow{(1_{DA}\Lotimes s)^{-1}}(DA)^{\Lotimes\ell+1}\xrightarrow{s\Lotimes1_{DA}}DA[m]\Big)\ \mbox{ in }\ \DDD(A^e).
\end{equation}
Since $\End_{\DDD(A^e)}(DA[m])\simeq\End_{\DDD(A^e)}(A)=Z(A)$, there is $z\in Z(A)^\times$ such that
\begin{equation}\label{define z}
t=(z\cdot)=(\text{multiplying }z).
\end{equation}
\new{We call $z$ the \emph{$\nu$-commutator} (of degree $\ell$), and denote by $z=z_A=z_A^{(\ell)}$.}\marginpar{A name is given. /OI}
Notice that $z$ does not depend on the choice of $s$ since a different choice of $s$ only differs by multiplying an element in $Z(A)^\times$.

\begin{theorem}\label{trivial extension is periodic}
Let $A$ be a finite dimensional algebra over a field $k$ such that $\gl A$ is finite and $A/{\rm rad A}$ is a separable $k$-algebra.
Assume that $A$ is $\frac{m}{\ell}$-Calabi--Yau, and define $z\in Z(A)^\times$ by \eqref{define z}.
Let $B=T(A)$ be the trivial extension algebra of $A$.
\begin{enumerate}[\rm(a)]
\item For the automorphism of $B$ given by $\varphi(a,f)=(a,(-1)^\ell zf)$ for $(a,f)\in A\oplus DA=B$, we have an isomorphism
\[\Omega^{\ell+m}_{B^e}(B)\simeq{}_\varphi B_1\ \mbox{ in }\mod B^e.\]
\item $z^\ell=1$ holds. In particular, for the order $q$ of $(-1)^\ell z$ in $A^\times$, we have
\[\Omega^{(\ell+m)q}_{B^e}(B)\simeq B\ \mbox{ in }\mod B^e.\]
\end{enumerate}
\end{theorem}

To prove this, let $P$ be a projective resolution of the $A^e$-module $DA$, where
\[P=[\cdots\xrightarrow{d^{-3}} P^{-2}\xrightarrow{d^{-2}} P^{-1}\xrightarrow{d^{-1}}P^0\to0\cdots].\]
The complex $C:=A\oplus P$ of $k$-vector spaces has a natural structure of a dg $k$-algebra such that $(a,f)\cdot (b,g)=(ab, ag+fb)$ for $(a,f),(b,g)\in A\oplus P=C$. Then we have a quasi-isomorphism
\[C\to H^0(C)=B\]
of dg $k$-algebras. We denote by $\CCC(C^e)$ the category of dg $C^e$-modules. For $i\ge0$, let
\begin{eqnarray}
P^{\otimes i}:=\overbrace{P\otimes_A\cdots\otimes_AP}^i\in\CCC(A^e),\\
Q_i:=C\otimes_AP^{\otimes i}\otimes_AC\in\CCC(C^e).
\end{eqnarray}
For $i >0$, define a morphism $f_i:Q_i\to Q_{i-1}$ in $\CCC(C^e)$ by
\begin{align*}
f_i(c_0\otimes c_1\otimes\cdots\otimes c_i\otimes c_{i+1})&=\sum_{j=0}^i(-1)^jc_0\otimes c_1\otimes\cdots c_jc_{j+1}\otimes\cdots\otimes c_i\otimes c_{i+1}\\
&=c_0c_1\otimes c_2\otimes\cdots\otimes c_{i+1}+(-1)^i
c_0\otimes \cdots\otimes c_{i-1}\otimes c_ic_{i+1},
\end{align*}
where $c_0,c_{i+1}\in C$ and other $c_j$'s are in $P$.
This, together with $f_0:Q_0\to C,\:c_0\otimes c_1\mapsto c_0c_1$, gives a bar-like resolution, that is, a complex
\[\cdots\xrightarrow{f_2}Q_1\xrightarrow{f_1}Q_0\xrightarrow{f_0}C\ \mbox{ in }\ \CCC(C^e)\]
whose total dg $C^e$-module is acyclic.
%It is classical and easy to show that the total complex of the bar resolution is acyclic.

We truncate this bar-like resolution using the dg $C^e$-module
\[N:=P^{\otimes\ell}\oplus P^{\otimes\ell+1},\]
whose $C^e$-action is given by $(x,y)(a,p):=(xa,ya+x\otimes p)$ and $(a,p)(x,y):=(ax,ay+p\otimes x)$ for $(x,y)\in P^{\otimes\ell}\oplus P^{\otimes\ell+1}=N$ and $(a,p)\in A\oplus P=C$.  We denote by
\[\sigma:C\to C\]
the automorphism of the dg $k$-algebra $C$ given by $\sigma(a,p)=(a,(-1)^\ell p)$ for $(a,p)\in A\oplus P=C$. We denote by ${}_\sigma N$ the dg $C^e$-module whose left action of $C$ is twisted by $\sigma$, that is, $(a,p)(x,y):=(ax,ay+(-1)^\ell p\otimes x)$.
We define a morphism
\[g_\ell:{}_\sigma N\to Q_{\ell-1}\ \mbox{ in }\ \CCC(A^e)\]
by $g_\ell(x,y)=x\otimes 1+(-1)^\ell 1\otimes x+y\in Q_{\ell-1}$.

The first key ingredient of the proof is the following.

\begin{lemma}\label{truncate bar resolution}
$g_\ell$ is a morphism in $\CCC(C^e)$ satisfying $f_{\ell-1}\circ g_\ell=0$. Moreover, the total dg $C^e$-module of
\begin{equation}\label{bar cpx}{}_\sigma N\xrightarrow{g_\ell}Q_{\ell-1}\xrightarrow{f_{\ell-1}}\cdots\xrightarrow{f_1}Q_0\xrightarrow{f_0}C\ \mbox{ in }\ \CCC(C^e)\end{equation}
is acyclic.
\end{lemma}

\begin{proof}
It is easy to check that $g_\ell$ is a morphism of graded $C^e$-modules, thanks to the twist $\sigma$. As morphisms in $\CCC(A^e)$, $g_\ell$ and $f_i$'s can be written as follows, where $\cdot:=\otimes_A$, $P^i:=P^{\otimes i}$ and $\epsilon:=(-1)^\ell$.
\[\xymatrix@R0em@C2em{
P^\ell\ar[dr]|(.35){\epsilon}\ar[ddr]|(.35)1&A\cdot P^{\ell-1}\cdot A\ar[dr]|(.65){-\epsilon}\ar[ddr]|(.65)1&\cdots\ar[dr]|(.35){-1}\ar[ddr]|(.35)1&A\cdot P^2\cdot A\ar[dr]|{1}\ar[ddr]|1&A\cdot P\cdot A\ar[dr]|{-1}\ar[ddr]|1&A\cdot A\ar[dr]|1\\
&A\cdot P^{\ell-1}\cdot P\ar[ddr]|(.65)1&\cdots\ar[ddr]|(.35)1&A\cdot P^2\cdot P\ar[ddr]|1&A\cdot P\cdot P\ar[ddr]|1&A\cdot P\ar[ddr]|1&A\\
P^{\ell+1}\ar[dr]|(.35)1&P\cdot P^{\ell-1}\cdot A\ar[dr]|(.65){-\epsilon}&\cdots\ar[dr]|(.35){-1}&P\cdot P^2\cdot A\ar[dr]|1&P\cdot P\cdot A\ar[dr]|{-1}&P\cdot A\ar[dr]|1\\
&P\cdot P^{\ell-1}\cdot P&\cdots&P\cdot P^2\cdot P&P\cdot P\cdot P&P\cdot P&P.}\] \marginpar{ED}
%
Therefore $g_\ell$ is a morphism in $\CCC(C^e)$ satisfying $f_{\ell-1}\circ g_\ell=0$, and the total dg module is contractible as a dg $A^e$-module, and hence acyclic.
\end{proof}

In particular, we obtain the following, where
\[\DDD^{\bo}(C^e):=\{X\in\DDD(C^e)\mid\dim H(X)<\infty\}\ \mbox{ and }\ \DDD_{\sg}(C^e):=\DDD^{\bo}(C^e)/\per C^e.\]

\begin{lemma}\label{C and N}
We have an isomorphism $C\simeq {}_\sigma N[\ell]$ in $\DDD_{\sg}(C^e)$.
\end{lemma}

\begin{proof}
Since $\gl A$ is finite and $A/{\rm rad} A$ is a separable $k$-algebra, it follows that $\gl A^e$ is also finite. Thus $P^{\otimes i}\in\per A^e$ and hence $Q_i\in\per C^e$ and $Q_i\simeq 0$ in $\DDD_{\sg}(C^e)$.

As usual, for a object $X,Y$ in a triangulated category $\TT$, we write
\[X*Y:=\{Z\in\TT\mid \mbox{there exists a triangle $X\to Z\to Y\to X[1]$ in $\TT$}\}.\]
The total dg module of \eqref{bar cpx} is in $C*Q_0[1]*\cdots*Q_{\ell-1}[\ell]*{}_\sigma N[\ell+1]$ and is isomorphic to the zero object in $\DDD(C^e)$ by Lemma~\ref{truncate bar resolution}.
Since each $Q_i\simeq 0$ in $\DDD_{\sg}(C^e)$, the zero object is contained in the subcategory $C*{}_\sigma N[\ell+1]$ of $\DDD_{\sg}(C^e)$. Thus there is a triangle $C\to 0\to{}_\sigma N[\ell+1]\to C[1]$ in $\DDD_{\sg}(C^e)$, and hence ${}_\sigma N[\ell]\simeq C$.
\end{proof}

The quasi-isomorphism $C\to H^0(C)=B$ gives a quasi-isomorphism $C^e\to H^0(C^e)=B^e$ and hence we have equivalences 
\[F:\DDD(C^e)\simeq\DDD(B^e)\ \mbox{ and }\ \DDD_{\sg}(C^e)
\simeq\DDD_{\sg}(B^e).\]
We denote by
\[\tau:B\to B\]
an automorphism of the $k$-algebra $B$ given by $\tau(a,f)=(a,zf)$ for $(a,f)\in A\oplus DA=B$.

\begin{lemma}\label{N and B}
We have $F(C)\simeq B$, and $F(N)\simeq{}_\tau B[m]$ in $\DDD(B^e)$.
\end{lemma}

To prove this, recall that, for each integer $\ell$, a full subcategory
\[\DDD^{\ell}(C^e):=\{X\in\DDD^{\bo}(C^e)\mid H^i(X)=0\mbox{ for all }i\neq\ell\}\subset\DDD(C^e)\]
is the heart of a shifted standard t-structure of $\DDD(C^e)$, and we have an equivalence (e.g.\ \cite[Proposition 4.8]{IY})
\begin{equation}\label{describe heart}
H^{\ell}:\DDD^{\ell}(C^e)\to\mod B^e.
\end{equation}

\begin{proof}[Proof of Lemma \ref{N and B}]
Clearly $F(C)\simeq B$ holds. 

We will show that $F(N)\simeq{}_\tau B[m]$. Since $s:(DA)^{\Lotimes\ell}\simeq A[m]$ in \eqref{define s} and $s\Lotimes1_{DA}:(DA)^{\Lotimes\ell+1}\simeq DA[m]$ are isomorphisms in $\DDD(A^e)$, we have $N\in\DDD^{-m}(C^e)$.
%We regard ${}_\tau B[m]$ as a dg $C^e$-module through the quasi-isomorphism $C^e\to B^e$. It suffices to show $N\simeq{}_\tau B$ in $\DDD(C^e)$.
Thanks to the equivalence \eqref{describe heart}, it suffices to show that $H^{-m}(N)$ is isomorphic to ${}_\tau B$ as $B^e$-module. 

Since $P$ is cofibrant as a dg $A^e$-module, the morphism $s$ in \eqref{define s} is given by a morphism
\[s:P^{\otimes\ell}\to A[m]\ \mbox{ in }\CCC(A^e).\]
Then $s\Lotimes1_{DA}$ is given by $s\otimes 1_P:P^{\otimes\ell+1}\to P[m]$. Consider the morphism
\[u:=s\oplus(s\otimes 1_P):N=P^{\otimes\ell}\oplus P^{\otimes\ell+1}\to C[m]=A[m]\oplus P[m]\ \mbox{ in }\ \CCC(A^e)\]
and the induced morphism
\[H^{-m}(u):H^{-m}(N)\to H^{-m}(C[m])=B\ \mbox{ in }\ \mod A^e.\]
Since $N$ is a dg $C^e$-module, $H^{-m}(N)$ has a natural $B^e$-module structure. To prove our claim, it suffices to show the following.
\begin{enumerate}
\item $H^{-m}(u)$ is a morphism of right $B$-modules.
\item If we twist the left action of $B$ on $B$ by $\tau$, then $H^{-m}(u)$ is a morphism of left $B$-modules.
\end{enumerate}
The claim (1) is clear. In fact, since $u=s\otimes1_C:N=P^{\otimes\ell}\otimes_C C\to C[m]$ is a morphism of right dg $C$-modules,  $H^{-m}(u)$ is a morphism of right $B$-modules.

It remains to show (2). It suffices to check that $H^{-m}(u)$ commutes with the left action of $DA\subset B$ after twisting by $\tau$, that is, for each $p\in Z^0(P)$, the diagram
\begin{equation}\label{show commutative}
\xymatrix{H^{-m}(P^{\otimes\ell})\ar[rr]^{H^{-m}(s)}\ar[d]_{H^{-m}(p\otimes-)}&& H^{-m}(A[m]) \ar@{=}[rr] &&A\ar[d]^{zH^0(p\cdot)}\\
H^{-m}(P^{\otimes\ell+1})\ar[rr]^{H^{-m}(s\otimes1_P)}&&H^{-m}(P[m])\ar[rr]^(.6)\sim&&DA}
\end{equation}
commutes, where the map $H^0(p\cdot)$ is induced from $(p\cdot):A\to P$.

By \eqref{define t}, we have a commutative diagram in $\DDD(A^e)$:
\[\xymatrix@R0.5em{
&&P[m]\ar[rr]^\sim&&DA[m]\ar[dd]^{t}\\
P^{\otimes\ell+1}\ar[urr]^{1_P\otimes s}\ar[drr]_{s\otimes1_P}\\
&&P[m]\ar[rr]^\sim&&DA[m].
}\]
Applying $H^{-m}$, we have a commutative diagram in $\mod A^e$:
\begin{equation}\label{know commutative}
\xymatrix@R0.5em{
&&H^{-m}(P[m])\ar[rr]^(.6)\sim&&DA\ar[dd]^{z\cdot}\\
H^{-m}(P^{\otimes\ell+1})\ar[urr]^(.47){H^{-m}(1_P\otimes s)}\ar[drr]_(.47){H^{-m}(s\otimes1_P)}\\
&&H^{-m}(P[m])\ar[rr]^(.6)\sim&&DA
}\end{equation}
since $t=(z\cdot)$ by \eqref{define z}. 
%\old{Therefore for each $p\in Z^0(P)$ and $x\in Z^{-m}(P^{\otimes\ell})$, applying the commutative diagram to $p\otimes x\in Z^{-m}(P^{\otimes\ell+1})$, we have
%\[z[ps(x)]=[(s\otimes 1_P)(p\otimes x)]\ \mbox{ in }DA,\]
%where $[y]\in H^{-m}(P[m])$ denote the cohomology class of $y\in P[m]$.
%Therefore the diagram \eqref{show commutative} commutes.}
Now fix $p\in Z^0(P)$ and $x\in Z^{-m}(P^{\otimes\ell})$.
For any $i\in \Z$, denote by $H^{-i}(y)$ the cohomology class of $y\in Z^{-i}(y)$.  Then sending $x$ through the upper composition in \eqref{show commutative} yields $zH^{0}(p)H^{-m}(s(x))=zH^{-m}(ps(x))$, which coincides with the image of $p\otimes x$ through the composition in the upper three maps of \eqref{know commutative}.
Likewise, the image of $p\otimes x$ through the lower composition in \eqref{know commutative} equals the image of $x$ under the composition of the three lower maps in \eqref{show commutative}.  Hence, the desired commutativity of \eqref{show commutative} follows from that of \eqref{know commutative}.
\end{proof}

Now we are ready to prove Theorem \ref{trivial extension is periodic}.

\begin{proof}[Proof of Theorem \ref{trivial extension is periodic}]
(a) In $\DDD_{\sg}(B^e)$, we have an isomorphisms
\[B\stackrel{\ref{N and B}}{\simeq}F(C)\stackrel{\ref{C and N}}{\simeq}F({}_\sigma N[\ell])\stackrel{\ref{N and B}}{\simeq}{}_{\sigma\tau}B[\ell+m].\]
Thus $\Omega^{\ell+m}_{B^e}(B)\simeq {}_{\sigma\tau}B$ holds in $\underline{\mod}B^e$. Thus the assertion follows.

(b) For simplicity, let $P^i:=(DA)^{\Lotimes i}$. Using \eqref{define t} twice, we have
\begin{align*}
&\Big(P^2[m]\xrightarrow{(1\otimes1\otimes s)^{-1}}P^{\ell+2}\xrightarrow{s\otimes1\otimes 1}P^2[m]\Big)\\
=&\Big(P^2[m]\xrightarrow{(1\otimes1\otimes s)^{-1}}P^{\ell+2}\xrightarrow{1\otimes s\otimes 1}P^2[m]\xrightarrow{(1\otimes s\otimes 1)^{-1}}P^{\ell+2}\xrightarrow{s\otimes1\otimes 1}P^2[m]\Big)=(z^2\cdot)\ \mbox{ in }\ \DDD(A^e).
\end{align*}
Repeating this, we have
\begin{equation}\label{z^l}
\Big(P^\ell[m]\xrightarrow{(1_{P^\ell}\otimes s)^{-1}}P^{2\ell}\xrightarrow{s\otimes1_{P^\ell}}P^\ell[m]\Big)=(z^\ell\cdot)\ \mbox{ in }\DDD(A^e).
\end{equation}
On the other hand, the diagram
\[\xymatrix{
P^{2\ell}\ar[d]_{1_{P^\ell}\otimes s}^\wr\ar[rr]^{s\otimes 
1_{P^\ell}}_\sim&&P^{\ell}[m]\ar[d]^{1_{A[m]}\otimes s}_\wr\\
P^{\ell}[m]\ar[rr]_{s\otimes 1_{A[m]}}^\sim&&A[2m]
}\]
commutes in $\DDD(A^e)$ since both compositions equal $s\otimes s$. Since $1_{A[m]}\otimes s=s[m]=s\otimes1_{A[m]}$ holds, we have $1_{P^\ell}\otimes s=s\otimes 1_{P^\ell}$.
%\[\Big(P^\ell[m]\xrightarrow{(s\otimes1^{\otimes\ell})^{-1}}A[2m]\xrightarrow{1^{\otimes\ell}\otimes s}P^\ell[m]\Big)=1_{P^\ell[m]}\ \mbox{ in }\ \DDD(A^e),\]}
By \eqref{z^l}, $(z^\ell\cdot)=1_{P^\ell[m]}$ holds. Thus $z^\ell=1$.
\end{proof}

\section{Calculation of $z_A$ in characteristic two and for Dynkin type path algebras}

\marginpar{The title is too special. We should give more systematic results and examples in this section. /OI}
For a finite dimensional algebra $A$ we denote by $A^\times$ the units of $A$ and by $Z(A)$ the center of $A$. Recall that finite dimensional algebras are Dedekind-finite so an element is a unit if and only if it has a left inverse (or a right inverse).\marginpar{Many claims should be written much earlier, e.g. 7.1 and 7.2 in Section 1.1. /OI}

\begin{lemma}
The units of the trivial extension algebra $T(A)$ of $A$ are given by $T(A)^\times = \{ (a,f) | a \in A^\times \}$.
The inverse of $(a,f) \in T(A)$ is given by $(a^{-1}, -a^{-1}f a^{-1})$.
\end{lemma}
\begin{proof}
The identity in $T(A)$ is given by $(1,0)$.
Let $(a,f) \in T(A)$ be a unit. Then there exists $(b,g) \in T(A)$ with
$(1,0)=(a,f)(b,g)=(ab,ag+fb)$. This implies that $a$ is invertible with $b=a^{-1}$. 
Now assume an element $(a,f) \in T(A)$ is given with $a$ invertible and define $(b,g)$ with $b:=a^{-1}$ and $g:=-a^{-1}f a^{-1}$.
Then $(a,f)(b,g)=(ab,ag+fb)=(1,0)$ and $(a,f)$ is invertible.
\end{proof}

\begin{proposition} \label{inneriffidentity}
Let $\phi_r: T(A) \rightarrow T(A)$ be the $K$-algebra automorphism of $T(A)$ given by $\phi_r(a,f)=(a,rf)$ for an element $r \in A^\times \cap Z(A)$.
Then $\phi_r$ is an inner automorphism if and only if $r=1$.
\end{proposition}

\begin{proof}
In case $r=1$, $\phi_r$ is the identity and thus inner.
Now assume that $\phi_r$ is inner. 
Thus for all $(a,f) \in T(A)$ we have for an invertible $(b,g) \in T(A)$:
$(a,rf)=(b,g)(a,f)(b,g)^{-1}=(ba,bf+ga)(b^{-1},-b^{-1}gb^{-1})=(bab^{-1},-bab^{-1}gb^{-1}+(bf+ga) b^{-1})$.
This implies that $a=bab^{-1}$ for all $a \in A$ and thus $b \in Z(A)$.
Then $-bab^{-1}gb^{-1}+(bf+ga) b^{-1}=-agb^{-1}+f+gab^{-1}=rf$ for all $(a,f) \in T(A)$.
Choosing $a=0$, we obtain $f=r f$ for all $f \in D(A)$, which implies $r=1$.

\end{proof}

For a given $\frac{m}{l}$-fractionally Calabi--Yau algebra $A$, we \new{consider the $\nu$-commutator $z=z_A^{(\ell)}\in  Z(A)^\times$ of degree $\ell$ in \eqref{define z}.} Note that $z_A$ also depends on the choosen $l$ and we sometimes denote it by $z_A^{(l)}$ when the dependency on $l$ is important. 

\begin{corollary} \label{periodcorollary}
Let $A$ be $\frac{m}{l}$-fractionally Calabi--Yau and assume $A$ has finite global dimension and $A/\rad A$ is a separable $k$-algebra.
Then $z_A=(-1)^l$ if and only if $B=T(A)$ has period dividing $l+m$.

\end{corollary}
\begin{proof}
Assume that $z_A=(-1)^l$. By \ref{trivial extension is periodic}, $\Omega_{B^e}^{l+m}(B) \cong _{\varphi}B_1$ in $\mod B^e$, where 
$\varphi(a,f)=(a,(-1)^l z_A f)=(a,f)$ is the identity and thus the period of $B$ divides $l+m$.
Now if the period of $B$ divides $l+m$, we have $B \cong \Omega_{B^e}^{l+m}(B) \cong _{\varphi}B_1$ and thus $\varphi$ is inner. This implies by \ref{inneriffidentity} that $\varphi$ is the identity and thus $z_A=(-1)^l$.
\end{proof}

The next theorem is the \new{derived version of} the well known Noether--Deuring theorem, see for example \cite[19.25]{Lam} \new{for modules version.}

\marginpar{Delete either 7.4 or 7.5. Written as Propositions since they are elementary. /OI}
\old{
\begin{proposition} \label{noetherdeuring}
Let $A$ be a finite dimensional $k$-algebra and $K$ a field extension of $k$. 
Then two finitely generated $A$-modules $M$ and $N$ are isomorphic as $A$-modules if and only if $K \otimes_k M$ and $K \otimes_k N$ are isomorphic as $K \otimes_k A$-modules.
\end{proposition}

We also need the following derived version of the Noether--Deuring theorem, that} We state here in the special case for finite dimensional algebras, see \cite[Theorem 3.8.8]{Z} for the general version.\marginpar{In \cite[3.8.8]{Z}, there is an extra assumption on the residue fields. /OI}
\begin{proposition} \label{noetherdeuringderived}
Let $A$ be a finite dimensional $k$-algebra and $K$ a field extension of $k$.
Let $X$ and $Y$ be two objects in $D^b(\mod A)$ such that $\End_{D^b(\mod A)}(X)$ is finite dimensional as a $k$-vector space. Then
$K \Lotimes_k X \cong K \Lotimes_k Y$ if and only if $X \cong Y$.
\end{proposition}

\begin{proposition} \label{characterstic2case}
Let $A$ be a $\frac{m}{l}$-fractionally Calabi--Yau algebra over the field $k$ and let $K$ be a field extension of $k$.
\begin{enumerate}
\item $A_K:=A \otimes_k K$ is also $\frac{m}{l}$-fractionally Calabi--Yau algebra for any field extension $K$ of $k$.
\item The periods of $T(A)$ and $T(A_K)$ coincide.
\item In case $k$ is the field with two elements and $K$ any field with characteristic two, we have $z_{A_K}=z_A=1$.
\end{enumerate}
\end{proposition}
\begin{proof}
\begin{enumerate}
\item This follows using \ref{noetherdeuringderived}.
\item Tensoring with $K$ a minimal projective bimodule resolution of $T(A)$ gives a minimal projective bimodule resolution of $T(A_K)$ and the period is preserved by \ref{noetherdeuring}.
\item Since $k$ has only one unique invertible element, $z_A=1$ and since $T(A)$ and $T(A_K)$ have the same period, \ref{periodcorollary} implies $z_{A_K}=1$.


\end{enumerate}

\end{proof}



\begin{proposition} \label{tensorproductforz}
Let $A$ be $\frac{m_1}{l_1}$-fractionally Calabi--Yau and let $B$ be $\frac{m_2}{l_2}$-fractionally Calabi--Yau.
Let $C:=A \otimes_K B$, which is $\frac{m}{l}$-fractionally Calabi--Yau by \ref{eg:fCYdim} with $l=\operatorname{lcm}(l_1,l_2)$ and $m= l (\frac{m_1}{l_1} + \frac{m_2}{l_2})$.
Then $z_{A \otimes_k B}^{(l)}= (z_A^{(l_1)})^{\frac{l}{l_1}} (z_B^{(l_2)})^{\frac{l}{l_2}}$.

\end{proposition}
\begin{proof}
TODO.
\end{proof}


We now calculate $z_A$ when $A$ is a path algebra of Dynkin type.
We first recall three results from the literature. Recall that the Coxeter number $h_A$ of a path algebra $A$ is equal to $n+1$ for type $A_n$, equal to $2(n-1)$ in type $D_n$, equal to 12 for type $E_6$, equal to 18 in type $E_7$ and equal to 30 in type $E_8$.


The next theorem is due to Brenner--Butler--King \cite[Theorem 2.1]{BBK}:

\begin{theorem} \label{BBKresult}
Let $k$ be a field of characteristic not equal two. Then $T(kQ)$ is periodic with period equal to $2(h-1)$.
\end{theorem}

For the next proposition, see \cite[Theorem 4.1]{MY}
\begin{theorem} \label{MYresult}
Let $A=kQ$ be a path algebra of Dynkin type. 
Then $A$ is $\frac{h_A-2}{h_A}$-fractionally Calabi--Yau.
\end{theorem}

\begin{proposition} \label{zcalculationhereditary}
Let $A=kQ$ be a path algebra of Dynkin type, which is $\frac{h_A-2}{h_A}$-Calabi--Yau.
Then we have for $z_{kQ}=(-1)^{n+1}$ when $Q$ is of Dynkin type $A_n$ and $z_{kQ}=1$ when $Q$ is not of Dynkin type $A_n$.

\end{proposition}
\begin{proof}
We can assume that $k$ has characteristic not equal to two, since the case of characteristic two is treated by \ref{characterstic2case} and the result is true in that case.
Using that $A$ is $\frac{h_A-2}{h_A}$-Calabi--Yau and that by \ref{BBKresult} $B=T(A)$ has period $2(h_A-1)$, we can apply \ref{periodcorollary} to conclude $z_A=(-1)^l=(-1)^{h_A}$, which proves the result since $h_A$ is even when $A$ is not of Dynkin type $A_n$.





\end{proof}

\new{
We give an example for the calculation of $z_A$ for a tensor product, which illustrates \ref{tensorproductforz}.
\begin{example}
Let $A=kQ_1$ and $B=kQ_2$ with $Q_1$ of linear oriented Dynkin type $A_2$ and $Q_2$ of linear oriented Dynkin type $A_3$.
$A$ is $\frac{1}{3}$-fractionally Calabi-Yau and $b$ is $\frac{2}{4}$-fractionally Calabi-Yau. Thus $C:= A \otimes_k B$ is $\frac{10}{12}$-fractionally Calabi-Yau and it is well-known that $C$ is derived equivalent to $k Q$ with $Q=E_6$, which implies that the trivial extension of $C$ has period 22 by \ref{BBKresult}.
Let $l=12$.
By \ref{trivial extension is periodic}, with $\varphi(a,f)=(a,(-1)^l z_C^{(l)} f)$ the automorphism of $C$, we have $\Omega_{C^e}^{22}(C) \cong _{\varphi}C_1$ in $\mod C^e$, which implies that $\varphi$ is inner and thus the identity. 
Thus $z_C^{(l)}=1$, which agrees with the result obtained from \ref{tensorproductforz} in this special case.
\end{example}}


%% \marginpar{in prelim}
%% \old{
%% The next proposition is due to Herschend-Iyama \cite[Proposition 1.4]{HerIya}
%% \begin{proposition} \label{HIresult}
%% Let $\Lambda_i$ be $\frac{m_i}{l_i}$-fractionally Calabi--Yau algebras for $i=1,...,t$.
%% Then the algebra $A= \bigotimes\limits_{i=1}^{t}{\Lambda_i}$ is $\frac{m}{l}$-fractionally Calabi--Yau with $l=lcm(l_1,...,l_t)$ and $m=l(\frac{m_1}{l_1}+\cdots \frac{m_t}{l_t})$.
%% \end{proposition}
%% }

\begin{corollary}
Let $B_t=\bigotimes\limits_{i=1}^{t}{k A_2}$ be the Boolean lattice for $t \geq 1$.
Then $T(B_t)$ is periodic with period equal to $t+3$ in case $k$ has characteristic 2 or $t$ is odd and period equal to $2(t+3)$ else.

\end{corollary}

\begin{proof}
Since $kA_2$ has Coxeter number $3$ and is $\frac{1}{3}$-fractionally Calabi--Yau, it follows from Example \ref{eg:fCYdim} that $B_t$ is $\frac{t}{3}$-fractionally Calabi--Yau. 
By \ref{tensorproductforz} we have $z_{kB_t} = z_{kA_2}^t=(-1)^{t}$.
Then $\varphi(a,f)=(a,(-1)^{3+t}f)$ for the map $\varphi$ of $kB_t$ as in \ref{trivial extension is periodic}.
Thus if $t$ is odd or the characteristic is two, $\varphi$ is the identity and the period divides $l+m=t+3$. If $t$ is even, $\varphi$ is not the identity and not inner by \ref{inneriffidentity} and thus the period divides $2(l+m)$ and is larger than $l+m$.
Now note that all simple $A$-modules have period equal to $l+m$ (TO SHOW?) and thus the period is divided by $l+m=3+t$, which proves the result.
\end{proof}

\new{
\marginpar{Aaron: make this new Sec 3, re-arrange?\\}
\marginpar{move part of text to intro}
We saw that being fractionally Calabi--Yau for an algebra $A$ of finite global dimension can be checked via the periodicity of the trivial extension of $A$. Many important classes of fractionally Calabi--Yau algebras are derived equivalent to incidence algebras of posets. For example any path algebra of Dynkin type is derived equivalent to the incidence algebra of a distributive lattice and recently it was shown that the incidence algebra of the Tamari lattice is fractionally Calabi--Yau in \cite{R} which answered a conjecture of Chapoton. Another open conjecture of Chapoton states that the distributive lattices of order ideals of the poset of positive roots of a finite root system is fractionally Calabi--Yau, see for example the introduction of \cite{Y}. The next proposition shows that the outer automorphism group for many posets is finite and thus the periodicity conjecture is true for trivial extensions of incidence algebras of such posets.
\begin{theorem}\cite[Corollary 7.3.7]{SO} \label{outerautoposet}
Let $P$ be a finite poset which contains an element $x \in P$ such that any other element in $P$ is comparable to $x$. Then any automorphism of the incidence algebra $A$ of $P$ is the composition of an inner automorphism of $A$ with an automorphism of $P$. In particular, the outer automorphism group of $P$ is finite.

\end{theorem}
This applies in particular to any bounded poset and thus to any lattice.


Our main result allows us to test whether a given incidence algebra of a poset is fractionally Calabi--Yau by looking at the periodicity of the trivial extension of the incidence algebra.
Since explicit quiver and relations are known for the trivial extension of incidence algebras by \cite[Corollary 3.12]{FP} we can check for periodicity in a quick way, for example using \cite{QPA}.

We give one non-trivial example of a well studied poset whose trivial extension is periodic, which has been found using \cite{QPA}.
\begin{example}
Let $L$ be the free distributive lattice on 3 generators, that is $L$ is the distributive lattice of order ideals of the Boolean lattice of a 3-set. 
Let $B$ be the trivial extension of the incidence algebra $A$ of $L$. Then every simple $B$-module $S$ is periodic and satisfies $\Omega^{14}(S)\cong S$. Since $L$ is a lattice, the periodicity conjecture holds for the trivial extension of $A$ by \ref{outerautoposet} and $B$ is periodic.
\end{example}


\begin{remark}
For the incidence algebra $A$ of a general distributive lattice $L$ we have \newline $\Ext_A^1(D(A),A) \neq 0$ since there exists indecomposable injective $A$-modules of projective dimension one, see \cite{IM}. Thus incidence algebras of distributive lattices never have $d$-cluster tilting modules for $d \geq 2$. We will see however later that in important cases, such as for Boolean lattices, trivial extensions of incidence algebras can have $d$-cluster tilting modules.
\end{remark}
}


%%%%%%%%%%%%%%%%%%%%%%%%%%%%%%%%%%%%%%%%%%%%%%%%%%%%%%%%%%%%%%%%%%%
%%%%%%%%%%%%%%%%%%%%%%%%%%%%%%%%%%%%%%%%%%%%%%%%%%%%%%%%%%%%%%%%%%%
%%%%%%%%%%%%%%%%%%%%%%%%%%%%%%%%%%%%%%%%%%%%%%%%%%%%%%%%%%%%%%%%%%%
%%%%%%%%%%%%%%%%%%%%%%%%%%%%%%%%%%%%%%%%%%%%%%%%%%%%%%%%%%%%%%%%%%%


\section{Further questions}

\begin{enumerate}
\item Are two fractionally CY algebras of finite global dimension derived equivalent iff their trivial extensions are derived equivalent? (One direction is known due to Rickard)
\item Are two fractionally CY algebras of finite global dimension derived equivalnet iff they have the same Coxeter polynomial? (This is probably false but would be very interesting even in special cases to check for derived equivalences)
\item Are the distributive lattices of root posets fractionally CY (conjecture of Chapoton).
\item Is there an intristic characterisation of periodic trivial extension algebras as a subset of periodic symmetric algebras? For representation-finite symmetric algebras, the trivial extensions $A$ are exactly those with $\Ext_A^1(M,M)=0$ for all indecomposable $A$-modules $M$.
\item Is the stable endomorphism ring of a d-cluster tilting module over a $d$-representation-finite algebra fractionally Calabi--Yau? \texttt{(Question of Osamu? Computer found no counterexamples yet.)}
\item Is the periodicity conjecture true for trivial extension algebras? (That is: When all simple modules are periodic, then the algebra is periodic)
\end{enumerate}

\begin{thebibliography}{10}

\bibitem[Am]{Am} C. Amiot, \emph{Cluster categories for algebras of global dimension 2 and quivers with potential}, Ann. Inst. Fourier (Grenoble) 59 (2009), no. 6, 2525--2590.

\bibitem[ASS]{ASS} I. Assem, D. Simson, A. Skowronski,  {\it Elements of the Representation Theory of Associative Algebra 1}, London Mathematical Society Student Texts, Volume 65, 2006. 

\bibitem[AS]{AS} I. Assem, A. Skowronski, {\it Iterated tilted algebras type $\widetilde{\mathbb{A}_n}$}, Math. Z., 195 (1987), 269--290

\bibitem[BBK]{BBK} S. Brenner, M. Butler, A. King, {\it Periodic Algebras which are Almost Koszul.} Algebras and Representation Theory, Volume 5, pages 331-368 (2002).
\bibitem[CIM]{CIM} A. Chan, O. Iyama, R. Marczinzik, \emph{Auslander-Gorenstein algebras from Serre-formal algebras via replication}, Adv. Math. 345 (2019), 222--262.

\bibitem[DI]{DI} E. Darp\"o, O. Iyama, \emph{$d$-representation-finite self-injective algebras}, Adv. Math. 362 (2020), 106932.

\bibitem[EH]{EH} K. Erdmann, T. Holm, \emph{Maximal $n$-orthogonal modules for selfinjective algebras}, Proc. Amer. Math. Soc. 136 (2008), no. 9, 3069--3078.

\bibitem[ES]{ES} K. Erdmann, A. Skowronski, {\it Periodic algebras.} In: Trends in Representation Theory of Algebras and Related Topics, EMS Series of Congress Reports , 2008.

\bibitem[FP]{FP} E. Fernandez, M. Platzeck, {\it Presentations of Trivial Extensions of Finite Dimensional Algebras and a Theorem of Sheila Brenner}, Journal of Algebra Volume 249, Issue 2, 2002, Pages 326-344.

\bibitem[FGR]{FGR} R. Fossum, P. Griffith, I. Reiten, {\it Trivial Extensions of Abelian Categories}, Lecture Notes in Mathematics Volume 456, Springer, 1975.
\bibitem[GR]{GR} P. Gabriel, A. Roiter, {\it Representations of finite-dimensional algebras}, Springer, 1992.

\bibitem[GSS]{GSS} E. L. Green, N. Snashall and {\O}. Solberg, \emph{The Hochschild cohomology ring of a self-injective algebra of finite representation type}, Proc. Amer. Math. Soc. 131 (2003), 3387–3393.

\bibitem[H]{H} N. Hanihara, \emph{Auslander correspondence for triangulated categories}, \new{Algebra Number Theory 14 (2020), no. 8, 2037--2058.}

\bibitem[Hap]{Hap} D. Happel, \emph{Triangulated categories in the representation theory of finite-dimensional algebras}, London Mathematical Society Lecture Note Series, 119. Cambridge University Press, Cambridge, 1988. 

\bibitem[HerIya]{HerIya}
{\sc M. Herschend, O. Iyama},
`{{$n$}-representation-finite algebras and twisted fractionally Calabi--Yau algebras}',
{\it Bull. London Math. Soc.} 43 (2011), no. 3, 449--466.

\new{\bibitem[HIMO]{HIMO} M. Herschend, O. Iyama, H. Minamoto, S. Oppermann, \emph{Representation theory of Geigle-Lenzing complete intersections}, to appear in Mem. Amer. Math. Soc., arXiv:1409.0668.}

\bibitem[HIO]{HIO}
{\sc M. Herschend, O. Iyama, S. Oppermann}, 
`{{$n$}-representation infinite algebras}',
{\it Adv. Math.} 252 (2014), 292--342.

\bibitem[HW]{HW} D. Hughes, J. Waschb\"usch, Trivial extensions of tilted algebras, Proc. London Math. Soc. (3) 46 (1983), no. 2, 347--364.

\new{\bibitem[HS]{HS} B. Huisgen-Zimmermann, M. Saorín, \emph{Geometry of chain complexes and outer automorphisms under derived equivalence}, Trans. Amer. Math. Soc. 353 (2001), no. 12, 4757--4777.}

\bibitem[IM]{IM} O. Iyama, R. Marczinzik, \emph{Distributive lattices and Auslander regular algebras} , https://arxiv.org/abs/2009.07170.

\bibitem[IY]{IY} O. Iyama, D. Yang, \emph{Silting reduction and Calabi--Yau reduction of triangulated categories}, Trans. Amer. Math. Soc. 370 (2018), no. 11, 7861--7898.

\new{
\bibitem[JL]{JL} C. Jensen, H. Lenzing, \emph{Model-theoretic algebra with particular emphasis on fields, rings, modules}, Algebra, Logic and Applications, 2. Gordon and Breach Science Publishers, New York, 1989.}

\bibitem[J]{J} H. Jin, \emph{Cohen-Macaulay differential graded modules and negative Calabi-Yau configurations}, to appear in Adv. Math.

\bibitem[K]{K} B. Keller, \emph{On triangulated orbit categories}, Doc. Math. 10 (2005), 551--581.

\bibitem[KY]{KY} S. Koenig, D. Yang, \emph{Silting objects, simple-minded collections, $t$-structures and co-$t$-structures for finite-dimensional algebras}, Doc. Math. 19 (2014), 403--438. 

\new{
\bibitem[MS]{MS} F. H. Membrillo-Hernández, L. Salmerón, \emph{A geometric approach to the finitistic dimension conjecture}, Arch. Math. (Basel) 67 (1996), no. 6, 448--456.
}

\bibitem[MY]{MY} J. Miyachi, A. Yekutieli, {\it Derived Picard Groups of Finite-Dimensional Hereditary Algebras.} Compositio Mathematica, Volume 129, Issue 3, 2001 , pp. 341-368. 

\bibitem[Lam]{Lam} T. Y. Lam, {\it A first course in noncommutative rings.} Graduate Texts in Mathematics 131, New York, 2001.

\bibitem[P]{P} M. Pressland, {\it Internally Calabi–Yau algebras and cluster-tilting objects}, Mathematische Zeitschrift volume 287, pages 555-585 (2017).
\bibitem[QPA]{QPA} The QPA-team, QPA - Quivers, path algebras and representations - a GAP package, Version 1.25; 2016 (https://folk.ntnu.no/oyvinso/QPA/).
\bibitem[R]{R} B. Rognerud, {\it The bounded derived categories of the Tamari lattices are fractionally Calabi--Yau}, https://arxiv.org/abs/1807.08503 

\bibitem[Ri]{Ri} C. Riedtmann, \emph{Representation-finite self-injective algebras of class $A\sb{n}$}, Representation theory, II (Proc. Second Internat. Conf., Carleton Univ., Ottawa, Ont., 1979), pp. 449--520, Lecture Notes in Math., 832, Springer, Berlin, 1980.

\bibitem[Sc]{Sc} S. Schroll, \emph{Trivial extensions of gentle algebras and Brauer graph algebras}, J. Algebra 444 (2015), 183--200.

\bibitem[Sk]{Sk} A. Skowronski, \emph{Selfinjective algebras: finite and tame type}, Trends in representation theory of algebras and related topics, 169--238, Contemp. Math., 406, Amer. Math. Soc., Providence, RI, 2006.

\bibitem[SY]{SY} A. Skowronski, K. Yamagata, \emph{Selfinjective algebras of quasitilted type}, Trends in representation theory of algebras and related topics, 639--708, EMS Ser. Congr. Rep., Eur. Math. Soc., Zürich, 2008.

%\bibitem[S]{S} D. Simson, {\it Linear representations of partially ordered sets and vector space categories}, Algebra, Logic, and Applications Vol 4, 1992.
\bibitem[SY]{SY} A. Skowronski, K. Yamagata, {\it Frobenius Algebras I: Basic Representation Theory}, EMS Textbooks in Mathematics, 2011.

\bibitem[SO]{SO} E. Spiegel, C. O'Donnell, {\it Incidence algebras}, Pure and Applied Mathematics, Marcel Dekker Inc., 1997.

\bibitem[T]{T} H. Tachikawa, \emph{Representations of trivial extensions of hereditary algebras}, Representation theory, II (Proc. Second Internat. Conf., Carleton Univ., Ottawa, Ont., 1979), pp. 579--599, Lecture Notes in Math., 832, Springer, Berlin, 1980.

\bibitem[V]{V} D. Voigt, \emph{Induzierte Darstellungen in der Theorie der endlichen, algebraischen Gruppen}, Mit einer englischen Einführung. Lecture Notes in Mathematics, Vol. 592. Springer-Verlag, Berlin-New York, 1977.

\bibitem[Y]{Y} E. Yildirim, {\it The Coxeter transformation on cominuscule posets}, Algebr. Represent. Theory 22 (2019), no. 3, 699-722. 
\bibitem[Z]{Z} A. Zimmermann, \emph{Representation theory. A homological algebra point of view.} Algebra and Applications, Springer 2014.
\end{thebibliography}
\end{document}
